% ============================================================================
%  BAO CAO (tieng Viet)
%  Khoi phuc nhan tu cap nhat cuc bo: co che suy giam cua cac phuong phap
%  dua tren hieu chuan phu tro, va mot uoc luong bat bien thang do.
%
%  Compile duoc bang CA HAI engine:
%     pdflatex VLI_baocao_vi.tex     (chay 3 lan)   -- hien thi dung
%     xelatex  VLI_baocao_vi.tex     (chay 3 lan)   -- NEN DUNG: copy-paste duoc
% ============================================================================
\documentclass[11pt]{article}

\usepackage{iftex}
\ifPDFTeX
  \usepackage[utf8]{inputenc}
  \usepackage[T5]{fontenc}
\else
  \usepackage{fontspec}   % mac dinh Latin Modern, khong can font he thong
\fi
\usepackage[vietnamese]{babel}

\usepackage[margin=2.2cm]{geometry}
\usepackage{amsmath,amssymb,amsthm,mathtools}
\usepackage{booktabs}
\usepackage{array}
\usepackage{pdflscape}
\usepackage[hidelinks]{hyperref}

% --- moi truong dinh ly, dung CHUNG mot bo dem de so thu tu lien tuc --------
\newtheorem{dinhly}{Định lý}
\newtheorem{bode}[dinhly]{Bổ đề}
\newtheorem{menhde}[dinhly]{Mệnh đề}
\newtheorem{hequa}[dinhly]{Hệ quả}
\newtheorem{giathiet}[dinhly]{Giả thiết}
\theoremstyle{remark}
\newtheorem{nhanxet}[dinhly]{Nhận xét}
\renewcommand{\proofname}{Chứng minh}

\newcommand{\R}{\mathbb{R}}
\newcommand{\E}{\mathbb{E}}
\newcommand{\one}{\mathbf{1}}
\newcommand{\pos}[1]{\left[#1\right]_+}
\newcommand{\simplex}{\Delta_n}
\DeclareMathOperator{\softmaxop}{softmax}
\DeclareMathOperator{\LSEop}{LSE}
\DeclareMathOperator*{\argmaxop}{arg\,max}
\DeclareMathOperator*{\argminop}{arg\,min}

\renewcommand{\arraystretch}{1.25}

\title{Khôi phục nhãn từ cập nhật cục bộ trong học liên kết:\\
cơ chế suy giảm của các phương pháp hiệu chuẩn bằng dữ liệu phụ trợ,\\
và một ước lượng bất biến thang độ}
\author{}
\date{\today}

\begin{document}
\maketitle

\begin{abstract}
Chúng tôi nghiên cứu bài toán khôi phục số lượng mẫu theo từng lớp trong một
cửa sổ cập nhật của một client, từ gradient tích lũy của bias ở lớp tuyến
tính cuối. Chúng tôi chứng minh hai đồng nhất thức chính xác liên hệ gradient
đó với \emph{khối lượng lỗi} theo lớp, rồi dùng chúng để phân loại chính xác
xem phương pháp nào suy giảm, ở mức nào, và vì sao. Cơ chế trung tâm: một ước
lượng kết hợp cập nhật quan sát được với một thống kê \emph{có độ lớn tuyệt
đối} đo từ bên ngoài cửa sổ bị tấn công thì tiến về một đại lượng
\emph{không phụ thuộc chút nào} vào nhãn cần tìm (Định lý~\ref{dl:supdo}), vì
chuẩn gradient bị chặn bởi lỗi huấn luyện (Bổ đề~\ref{bd:chuan}) còn thống kê
kia thì không. Chúng tôi chứng minh ZLG$^+$ và bước đầu của RLU \emph{sụp đổ}
theo cơ chế này, với một biên tường minh cho đầu ra nguyên
(Hệ quả~\ref{hq:apdung}); rằng iLRG thất bại vì một lí do khác hẳn --- điều
kiện số (Mục~\ref{sec:ilrg}); và rằng đầu ra của LLG$^+$ cuối cùng chỉ còn
phụ thuộc nhãn qua \emph{mẫu dấu} của tín hiệu, tức không quá $C$ bit
(Mệnh đề~\ref{md:llgp}).
Sau đó chúng tôi trình bày VLI và dạng đóng VLI$^*$, với tính chất định danh
là bất biến theo thang độ của cập nhật (Mệnh đề~\ref{md:batbien}); dưới một
giả thiết đối xứng mà chúng tôi chứng minh là lựa chọn \emph{duy nhất} tương
thích với tính hoán vị nhãn (Mệnh đề~\ref{md:hoanvi}), VLI$^*$ thu về một
biểu thức không thang độ duy nhất và \emph{chính xác}
(Định lý~\ref{dl:dangdong}). Chúng tôi phát biểu các chế độ thất bại của giả
thiết đó một cách định lượng, thay vì khẳng định rằng nó đúng.
\end{abstract}

\tableofcontents

% ============================================================================
\section{Phạm vi và đóng góp}

\paragraph{Những gì báo cáo này thiết lập.}
\begin{enumerate}
\item Hai đồng nhất thức chính xác cho gradient bias của lớp cuối
  (Bổ đề~\ref{bd:chuan} và~\ref{bd:dongnhat}), bao gồm tính chất tổng bằng
  không $\sum_j G_j = 0$ --- tính chất này về sau làm phần lớn công việc.
\item Một định lý (Định lý~\ref{dl:supdo}) tách làm hai phần --- định tính
  chỉ cần \emph{tính liên tục tại một điểm}, định lượng cần Lipschitz
  \emph{cục bộ} --- cùng Bổ đề~\ref{bd:lipcucbo} cung cấp hằng số Lipschitz
  cục bộ tường minh $2n/\delta$. Nhờ đó Hệ quả~\ref{hq:apdung} \emph{chứng
  minh được} rằng ZLG$^+$ và RLU thỏa cả hai phần, điều mà một phát biểu đòi
  Lipschitz toàn cục không làm được (Mệnh đề~\ref{md:lipschitz}).
\item Một phân loại chính xác thay cho khẳng định ``mọi phương pháp đều sụp'':
  ZLG$^+$ và RLU sụp về một dự đoán độc lập nhãn; iLRG thất bại do điều kiện
  số, không do hiệu chuẩn; đầu ra của LLG$^+$ chỉ còn phụ thuộc nhãn qua mẫu
  dấu $h_1$, nên mang không quá $C$ bit (Mệnh đề~\ref{md:llgp}). LLG thuần thì
  bất biến thang độ ở cả hai pha, điều này giải thích vì sao nó bền vững.
\item Một chỉnh lại về nhân quả (Mục~\ref{sec:tuongtac}): việc tín hiệu co nhỏ
  \emph{không} tự mình phá vỡ các tấn công --- thông tin nhãn vẫn khôi phục
  được chính xác từ đại lượng thực mà ZLG$^+$ tính ra trước khi lượng tử hóa
  (Mệnh đề~\ref{md:conthongtin}). Thất bại là \emph{tương tác} giữa một
  \emph{lệch thang độ} $\|c\|\gg\beta\|\Gamma\|$ và một thiết kế so độ lớn
  tuyệt đối. Lệch thang độ có hai nguồn: với ZLG$^+$ và RLU nó là cấu trúc
  ($c$ có cỡ một vector xác suất, bất kể khoảng cách tổng quát hóa); với
  LLG$^+$ nó là khoảng cách tổng quát hóa. Cách sửa vì vậy là bất biến thang
  độ chứ không phải ước lượng phụ trợ tốt hơn.
\item Hai mệnh đề bổ trợ định lượng hai khiếm khuyết riêng của RLU: khoảng
  cách Jensen $e^{\sigma^2/2}$ khi lấy softmax của logit trung bình
  (Mệnh đề~\ref{md:jensen}), và việc phép cắt-rồi-chuẩn-hóa khuếch đại nhiễu
  (Mệnh đề~\ref{md:lipschitz}).
\item Một tính chất bất biến thang độ (Mệnh đề~\ref{md:batbien}) mà ước lượng
  đề xuất thỏa mãn còn các phương pháp cũ thì không, cùng một dạng đóng
  (Định lý~\ref{dl:dangdong}) chính xác dưới giả thiết đối xứng đã nêu.
\item Một lập luận về tính duy nhất (Mệnh đề~\ref{md:hoanvi}) thay thế cách
  biện minh thông thường ``cho đơn giản'' của giả thiết đối xứng bằng ``lựa
  chọn duy nhất khả dụng đối với một kẻ tấn công không có tri thức riêng về
  lớp''.
\item Một điều kiện đủ \emph{yếu hơn hẳn} giả thiết đó
  (Mệnh đề~\ref{md:dieukienyeu}), mở rộng vùng chính xác của dạng đóng.
\end{enumerate}

\paragraph{Những gì báo cáo này \emph{không} thiết lập.}
Vector độ khó $r$ và ma trận nhầm lẫn $Q$ giới thiệu dưới đây \emph{không}
được khẳng định là khôi phục được từ cập nhật được chia sẻ; ước lượng tổng
quát của Định lý~\ref{dl:tongquat} nhận chúng làm đầu vào và ở đây được dùng
như một \emph{trần lý thuyết} (oracle). Các bảng số ở
Mục~\ref{sec:kiemchung} là đánh giá trên mô hình sinh do các giả thiết hàm ý,
\textbf{không phải thí nghiệm trên một mạng đã huấn luyện}; chúng kiểm chứng
phần đại số và đo độ nhạy, và được ghi rõ như vậy ở mọi nơi.

% ============================================================================
\section{Kí hiệu}
\label{sec:kihieu}

Một cửa sổ của client chứa $n$ mẫu thuộc $C$ lớp. Với mẫu thứ $i$:

\begin{center}
\footnotesize
\begin{tabular}{@{}ll@{}}
\toprule
Kí hiệu & Ý nghĩa \\
\midrule
$y_i \in \{1,\dots,C\}$ & nhãn thật của mẫu $i$ \\
$z_i \in \R^C$ & vector logit, $z_i = W e_i + b$ với $e_i$ là embedding \\
$p_i = \softmaxop(z_i)$ & phân phối lớp mà mô hình dự đoán \\
$e_{y_i}$ & vector one-hot của lớp đúng \\
$\ell_i = -\log p_{i,y_i}$ & hàm mất mát cross-entropy \\
$L_{\mathrm{tr}} = n^{-1}\sum_i \ell_i$ & mất mát huấn luyện trung bình trên cửa sổ \\
\midrule
$N_j = \#\{i : y_i = j\}$ & số mẫu thật của lớp $j$; $\sum_j N_j = n$ \\
$\pi_j = N_j/n$ & tỉ lệ lớp thật \\
$P_j = \sum_i p_{i,j}$ & tổng khối lượng dự đoán dồn vào lớp $j$ \\
$G = \sum_{i=1}^n (p_i - e_{y_i})$ & \textbf{tín hiệu gradient bias} \\
\midrule
$\epsilon_i = 1 - p_{i,y_i}$ & lỗi của mẫu $i$ \\
$\bar\epsilon = n^{-1}\sum_i \epsilon_i$ & lỗi trung bình trên cửa sổ \\
$\bar\epsilon_j = N_j^{-1}\sum_{i:y_i=j}\epsilon_i$ & lỗi trung bình của lớp $j$ (lớp có mặt) \\
$a_j = N_j \bar\epsilon_j$ & \textbf{khối lượng lỗi} của lớp $j$; $a_j = 0$ nếu lớp vắng mặt \\
$r_j = \bar\epsilon_j/\bar\epsilon$ & độ khó tương đối; $a_j = N_j r_j \bar\epsilon$ \\
$\kappa = r_{\max}/r_{\min}$ & độ phân tán độ khó trên các lớp có mặt \\
\midrule
$E_j = \sum_{i:y_i\ne j} p_{i,j}$ & khối lượng lỗi chảy \emph{vào} lớp $j$ \\
$Q_{jk}$ & tỉ lệ khối lượng lỗi của lớp $j$ chuyển sang lớp $k$; $Q_{jj}=0$ \\
$v$ & vector Perron của $Q^T$: $Q^Tv = v$, $\one^Tv=1$ \\
\midrule
$M_i = z_{i,y_i} - \LSEop_{k\ne y_i} z_{i,k}$ & \emph{margin}; $\epsilon_i = (1+e^{M_i})^{-1}$ \\
$\simplex = \{x \in \R^C : x \ge 0,\ \one^Tx = n\}$ & simplex xác suất được phóng lên $n$ \\
$\pos{x} = \max(x,0)$ & phần dương theo từng thành phần \\
\bottomrule
\end{tabular}
\end{center}

Ma trận chuyển lỗi ngoài đường chéo được định nghĩa khi $a_j > 0$ bởi
\[
 Q_{jk} = \frac{\sum_{i : y_i = j} p_{i,k}}{a_j} \quad (k \ne j),
 \qquad Q_{jj} = 0 .
\]
Mỗi hàng được định nghĩa đều không âm và có tổng bằng một, vì
$\sum_{k\ne j}\sum_{i:y_i=j} p_{i,k} = \sum_{i:y_i=j}(1-p_{i,j}) = a_j$. Một
lớp có khối lượng lỗi bằng không có thể được gán phân phối tùy ý trên các lớp
còn lại; nó không đóng góp gì vào $Q^Ta$.

\paragraph{Điểm trục của toàn bộ báo cáo.}
Phép chia cho $a_j$ chính là thứ \emph{loại bỏ} $\bar\epsilon$ khỏi $Q$: ma
trận này mang \emph{hình dạng} nhầm lẫn của mô hình, không mang \emph{độ lớn}
của nó. Tương tự với $r$, vốn là một tỉ số giữa hai lỗi. Sự phân biệt giữa
\emph{hình dạng} và \emph{độ lớn} là bản lề của mọi lập luận phía sau.

\paragraph{Quy ước về tín hiệu.}
Một cài đặt thực tế tích lũy gradient của mất mát \emph{trung bình} theo từng
minibatch. Khi mọi minibatch có kích thước $B$ và mỗi cửa sổ gồm $m$ bước
($n = mB$), tín hiệu tích lũy sai khác với tổng theo mẫu $G$ ở trên bởi một
hệ số dương chung, và có thể thêm một phép chuẩn hóa cửa sổ chung nữa. Mọi
công thức được tuyên bố là bất biến thang độ dưới đây đều không bị ảnh hưởng
bởi hệ số như vậy. Dưới SGD thuần, cập nhật trọng số ở server có dấu ngược
lại; ở đây $G$ luôn kí hiệu gradient của mất mát.

% ============================================================================
\section{Các đồng nhất thức chính xác}

\begin{bode}[Chuẩn gradient bị kiểm soát bởi lỗi huấn luyện]
\label{bd:chuan}
Với mọi mẫu, $\|p_i - e_{y_i}\|_1 = 2\epsilon_i$. Do đó
\[
 \|G\|_1 \;\le\; 2n\bar\epsilon \;\le\; 2n L_{\mathrm{tr}} .
\]
\end{bode}

\begin{proof}
Thành phần ứng với lớp đúng của $p_i - e_{y_i}$ có giá trị tuyệt đối
$1 - p_{i,y_i} = \epsilon_i$. Các thành phần còn lại đều không âm và có tổng
bằng $1 - p_{i,y_i} = \epsilon_i$. Cộng lại được $2\epsilon_i$. Bất đẳng thức
tam giác cho $\|G\|_1 \le \sum_i \|p_i - e_{y_i}\|_1 = 2n\bar\epsilon$. Cuối
cùng $1 - x \le -\log x$ với $x \in (0,1]$; áp dụng tại $x = p_{i,y_i}$ rồi
lấy trung bình.
\end{proof}

Bổ đề~\ref{bd:chuan} là phát biểu định lượng của hiện tượng mà báo cáo này
giải thích: \emph{khi mô hình khớp dần dữ liệu huấn luyện, toàn bộ tín hiệu
khả dụng với bất kì kẻ tấn công nào đều co về không, tỉ lệ với lỗi huấn
luyện.} Nó \emph{không} nói rằng thông tin nhãn biến mất ---
Bổ đề~\ref{bd:dongnhat} cho thấy thông tin còn nguyên trong \emph{hướng} của
$G$ --- nhưng nó nói rằng mọi quy trình đem $G$ so với một độ lớn ngoại lai
cố định thì cuối cùng sẽ bị đại lượng ngoại lai đó chi phối.

\begin{bode}[Các đồng nhất thức về số lượng, khối lượng lỗi và nhầm lẫn]
\label{bd:dongnhat}
Đặt $P_j = \sum_i p_{i,j}$. Khi đó
\[
 G_j = P_j - N_j = -a_j + E_j,
 \qquad N_j = P_j - G_j,
 \qquad E_j = (Q^Ta)_j ,
\]
và do đó
\begin{equation}
 G = -(I - Q^T)\,a .
 \label{eq:hept}
\end{equation}
Hơn nữa
\begin{equation}
 \sum_{j=1}^C G_j = 0 ,
 \label{eq:tongkhong}
\end{equation}
và nếu $N_j = 0$ thì $G_j = P_j \ge 0$. Cuối cùng, vì
$a_j = N_j r_j \bar\epsilon$, tín hiệu \emph{phân tích thành nhân tử}
\begin{equation}
 G = \bar\epsilon \cdot \Bigl[-(I-Q^T)\,(N \odot r)\Bigr],
 \label{eq:nhantu}
\end{equation}
trong đó $\odot$ là phép nhân theo từng thành phần, và theo
Bổ đề~\ref{bd:chuan} vector trong ngoặc có $\|\cdot\|_1 \le 2n$.
\end{bode}

\begin{proof}
Thành phần thứ $j$ của $G$ cộng dồn toàn bộ khối lượng dự đoán dồn vào lớp
$j$ rồi trừ một đơn vị cho mỗi mẫu thuộc lớp $j$, cho ra $P_j - N_j$. Tách
$P_j$ thành phần của chính lớp đúng
$\sum_{i:y_i=j}p_{i,j} = N_j - a_j$ và phần của các lớp khác $E_j$; ta được
$G_j = -a_j + E_j$ và tính chất về dấu ở lớp vắng mặt. Với dạng nhầm lẫn, với
mỗi lớp đích $k$,
\[
 E_k = \sum_{j \ne k}\ \sum_{i : y_i = j} p_{i,k}
     = \sum_j a_j Q_{jk} = (Q^Ta)_k ,
\]
thế vào $G = -a + E$ ta được~\eqref{eq:hept}. Với~\eqref{eq:tongkhong}, ta có
$\sum_j P_j = \sum_i \sum_j p_{i,j} = n$ và $\sum_j N_j = n$.
Đẳng thức~\eqref{eq:nhantu} là việc thay $a = \bar\epsilon\,(N\odot r)$
vào~\eqref{eq:hept} rồi rút $\bar\epsilon$ ra ngoài.
\end{proof}

\begin{nhanxet}[Phân tích nhân tử nói gì và không nói gì]
Đẳng thức~\eqref{eq:nhantu} là chính xác, nhưng vector trong ngoặc phụ thuộc
vào $Q$ và $r$ của cửa sổ, mà hai đại lượng này \emph{thay đổi} khi mô hình
huấn luyện (Mục~\ref{sec:kiemchung} đo được $Q$ trôi tới $151\%$). Vì vậy
không được đọc~\eqref{eq:nhantu} là ``$G$ chỉ đổi thang độ khi $\bar\epsilon$
giảm''. Điều nó cho thấy là: $\bar\epsilon$ là một nhân tử chung, còn nhân
tử còn lại bị chặn ($\|\cdot\|_1\le 2n$), và có thể bằng không trong trường
hợp suy biến (cửa sổ cân bằng với $Q$ đều cho $G \equiv 0$).
\end{nhanxet}

\begin{nhanxet}
Đồng nhất thức~\eqref{eq:tongkhong} có ba hệ quả được dùng lại nhiều lần.
(i) Một $G$ khác không \emph{buộc phải} có một thành phần dương, nên cái mốc
mà VLI$^*$ sử dụng luôn tồn tại miễn tín hiệu không đồng nhất bằng không.
(ii) Việc bảo toàn khối lượng trong Định lý~\ref{dl:dangdong} là tự động, nên
\emph{không cần} bước chuẩn hóa nào. (iii) Bản thân gradient trung bình của
mỗi minibatch cũng có tổng bằng không, vì $B^{-1}(B - B) = 0$; tín hiệu tích
lũy do đó thừa hưởng~\eqref{eq:tongkhong} tới sai số dấu phẩy động, dù mô
hình có dịch chuyển giữa các bước.
\end{nhanxet}

\begin{nhanxet}[Tính đúng đắn qua nhiều bước cục bộ]
Bổ đề~\ref{bd:chuan} và~\ref{bd:dongnhat} là các phát biểu theo từng mẫu và
vẫn đúng khi các mẫu khác nhau được đánh giá tại các trạng thái mô hình cục
bộ khác nhau. Chúng \emph{không} tự động mô tả cập nhật của optimizer ở server
khi cập nhật đó còn chứa trạng thái momentum, weight decay, số hạng proximal,
hay số hạng chỉnh của FedDyn. Mục~\ref{sec:gioihan} sẽ quay lại điểm này.
\end{nhanxet}

% ============================================================================
\section{Các phương pháp cơ sở}
\label{sec:coso}

Chúng tôi viết lại mỗi phương pháp đã công bố đúng theo dạng nó được cài đặt,
để mục tiếp theo có thể chỉ ra chính xác đại lượng nào được mỗi phương pháp
vay từ bên ngoài cửa sổ bị tấn công. Các đại lượng có gạch ngang hay mũ
($\bar p$, $\bar O$, $\mathrm{shift}$, $\mathrm{offset}$) được đo trên một
tập \emph{phụ trợ} (auxiliary) mà kẻ tấn công có --- trong thực tế là một
phần được phân tầng của tập kiểm tra, nên xấp xỉ cân bằng theo lớp.

\paragraph{Hai tín hiệu mà mã nguồn thật sự dùng.}
Mã nguồn tích lũy gradient của mất mát \emph{trung bình} theo batch rồi chia
cho $m$, nên tín hiệu bias mà RLU và VLI$^*$ nhận là $G/n$. Còn LLG,
LLG$^+$ và ZLG$^+$ không dùng gradient bias mà dùng gradient trọng số cộng
theo chiều embedding, $u_j = \sum_d \Delta W_{jd}$. Vì
$u_j = n^{-1}\sum_i (p_{i,j}-\mathbb{1}[y_i=j])\,s_i$ với
$s_i = \sum_d e_{i,d}$, ta viết $u = \bar O\,\tilde G/n$, trong đó
\[
 \tilde G_j = \sum_i w_i\,(p_{i,j} - \mathbb{1}[y_i=j]),
 \qquad w_i = s_i/\bar O .
\]
Với embedding sau ReLU thì $s_i \ge 0$, nên $\tilde G$ vẫn có tổng bằng không
(như~\eqref{eq:tongkhong}) và $\|\tilde G\|_1 \le 2n\,w_{\max}\,\bar\epsilon$
(như Bổ đề~\ref{bd:chuan}, thêm nhân tử $w_{\max} = \max_i w_i$). Mọi kết
luận ``sụp đổ'' cho ba phương pháp này do đó cần thêm giả thiết $w_{\max}$
bị chặn, tức không có mẫu nào có tổng embedding lớn bất thường so với trung
bình phụ trợ. Khi mọi $w_i = 1$ thì $\tilde G = G$. Dưới đây $\Gamma$ kí
hiệu chung tín hiệu mà một phương pháp dùng ($G$ hoặc $\tilde G$).

\paragraph{iLRG (Ma và cộng sự, 2023).}
Khôi phục ước lượng embedding theo lớp bằng tỉ số giữa gradient trọng số và
gradient bias, $\hat e_j = \nabla_{W_j}\!/\,\nabla_{b_j}$, đưa nó qua bộ phân
lớp để thu các vector xác suất khôi phục, rồi giải hệ tuyến tính tìm số lượng
bằng giả nghịch đảo.

\paragraph{LLG (Wainakh và cộng sự, 2021).}
Với tín hiệu $u = \sum_d \Delta W$, tính một đại lượng vô hướng
\emph{impact} từ chính tín hiệu quan sát được,
\[
 \mathrm{impact} = \Bigl(1 + \tfrac1C\Bigr)\frac{\sum_{j : u_j<0} u_j}{B}
 \;<\; 0,
\]
chọn mọi lớp có thành phần âm (tập $h_1$), rồi lần lượt gán các suất còn
lại cho thành phần đang là $\argminop$; mỗi lần gán, thành phần đó được
\emph{cộng thêm} $|\mathrm{impact}|$ (mã nguồn viết \texttt{add(-impact)}).

\paragraph{LLG$^+$ (Wainakh và cộng sự, 2021).}
Như LLG, nhưng $h_1$ được chọn trên $u$ còn vòng tham lam chạy trên
$u - \mathrm{offset}$, và impact được đo trên dữ liệu phụ trợ:
\[
 \mathrm{offset}_j = \text{trung bình thành phần $j$ khi nhãn thật $\ne j$},
 \qquad
 \mathrm{impact} \approx -\Bigl(1+\tfrac1C\Bigr)\frac{\bar O\,\bar\epsilon_{\mathrm{aux}}}{B}
\]
(lấy trung bình với giá trị đo trước khi huấn luyện). Cả hai đều có độ lớn
tỉ lệ với $\bar\epsilon_{\mathrm{aux}}$.

\paragraph{ZLG$^+$ (Geng và cộng sự, 2021).}
Với $u = \bar O\,\tilde G/n$ như trên,
\begin{equation}
 \widehat N \;\propto\; \Bigl[n\Bigl(\bar p - \frac{u}{\bar O}\Bigr)\Bigr]_+
 \;=\; \bigl[n\bar p - \tilde G\bigr]_+ ,
 \label{eq:zlg}
\end{equation}
rồi nhân lại cho tổng bằng $n$; mã nguồn có thêm \texttt{int} trước và
\texttt{round} sau phép nhân lại (Hệ quả~\ref{hq:apdung}). Ở đây $\bar p_j$
là độ tin cậy trung bình vào lớp đúng $j$ trên dữ liệu phụ trợ, và $\bar O$ là
tổng embedding trung bình trên dữ liệu phụ trợ.

\emph{Ghi chú nguồn gốc:} việc ước lượng $\bar p$, $\bar O$ từ dữ liệu phụ
trợ \textbf{thật} không phải thiết kế gốc của Geng và cộng sự --- bài báo
gốc ước lượng cả hai đại lượng này bằng cách cho \emph{dữ liệu giả} (ảnh
sinh tổng hợp) đi qua mô hình (một cấu trúc zero-shot, đúng với động cơ họ
tuyên bố: không cần tập dữ liệu ngoài nào). Dạng dùng dữ liệu phụ trợ thật
ở đây và trong suốt báo cáo này theo đúng baseline ``ZLG$^+$'' mà Chen và
Vikalo (RLU) đặt ra để so sánh công bằng giữa các phương pháp, tất cả đều
dùng chung tập phụ trợ $A$; bản dữ liệu giả gốc được gọi là ZLG (không có
dấu $+$) ở những chỗ cần phân biệt hai dạng.

\paragraph{RLU (Chen và Vikalo, 2024).}
Với tín hiệu bias trung bình $G/n$ như mã nguồn,
\begin{equation}
 \widehat N = n \cdot \frac{\pos{\mathrm{shift}\cdot\bar\rho - G/n}}
                           {\bigl\|\pos{\mathrm{shift}\cdot\bar\rho - G/n}\bigr\|_1},
 \qquad \mathrm{shift} = \softmaxop(\mu_{\mathrm{aux}}),
 \label{eq:rlu}
\end{equation}
trong đó $\mu_{\mathrm{aux}}$ là vector \emph{logit trung bình} theo lớp đo
trên dữ liệu phụ trợ và $\bar\rho$ là một hệ số vô hướng tích lũy theo bước.
Đây là \emph{lần gọi đầu}; sau đó một thủ tục tìm kiếm hậu nghiệm điều chỉnh
$\mathrm{shift}$ và gọi lại~\eqref{eq:rlu} lần nữa.

\begin{center}
\footnotesize
\begin{tabular}{@{}llll@{}}
\toprule
Phương pháp & Đại lượng vay từ ngoài & Đơn vị & Dạng \\
\midrule
iLRG      & không (dùng $\nabla W/\nabla b$) & --- & giả nghịch đảo \\
LLG       & không                            & --- & tự hiệu chuẩn, tham lam \\
LLG$^+$   & $\mathrm{offset}$, impact        & tuyệt đối, $\propto \bar\epsilon_{\mathrm{aux}}$ & cộng \\
ZLG$^+$   & $\bar p$, $\bar O$               & tuyệt đối, $\bar p \approx 1$                    & \eqref{eq:zlg} \\
RLU       & $\mathrm{shift}$                 & tuyệt đối, vector xác suất ($\|\cdot\|_1=1$)    & \eqref{eq:rlu} \\
\bottomrule
\end{tabular}
\end{center}

% ============================================================================
\section{Cơ chế thất bại}
\label{sec:thatbai}

\subsection{Hiệu chuẩn chỉ bằng dữ liệu phụ trợ trở thành độc lập với nhãn}

\begin{dinhly}[Sự sụp đổ của các ước lượng hiệu chuẩn bằng dữ liệu phụ trợ]
\label{dl:supdo}
Giả sử một ước lượng tỉ lệ lớp có dạng
\[
 \widehat\pi = \Phi\bigl(c - \beta G\bigr),
\]
trong đó $c$ và hệ số $\beta>0$ chỉ phụ thuộc vào thông tin bên ngoài cửa sổ
bị tấn công. Khi đó:
\begin{enumerate}
\item[(i)] \textup{(Định tính.)} Nếu $\Phi$ \emph{liên tục tại} $c$ thì
 $\widehat\pi \to \Phi(c)$ khi $\bar\epsilon \to 0$ --- một vector được xác
 định hoàn toàn bởi $c$ và \emph{không} bởi các nhãn cần tìm.
\item[(ii)] \textup{(Định lượng.)} Nếu thêm nữa $\Phi$ là $L_\Phi$-Lipschitz
 trên một lân cận của $c$ chứa mọi $c-\beta G$ đang xét, thì
\begin{equation}
 \bigl\|\widehat\pi - \Phi(c)\bigr\|
 \;\le\; L_\Phi\,\beta\,\|G\|_1
 \;\le\; 2 L_\Phi\,\beta\, n\,\bar\epsilon
 \;\le\; 2 L_\Phi\,\beta\, n\, L_{\mathrm{tr}} .
 \label{eq:supdo}
\end{equation}
\end{enumerate}
Ở đây $\Phi$ có thể trả về tỉ lệ hoặc số lượng; nếu nó trả về số lượng (như
$R$ ở Bổ đề~\ref{bd:lipcucbo}) thì thừa số $n$ đã nằm sẵn trong $L_\Phi$.
Phát biểu vẫn đúng khi thay $G$ bởi tín hiệu có trọng số $\tilde G$, với
$\|\tilde G\|_1 \le 2nw_{\max}\bar\epsilon$ (Mục~\ref{sec:coso}).
\end{dinhly}

\begin{proof}
(i) Theo Bổ đề~\ref{bd:chuan}, $\|\beta G\|_1 \le 2\beta n\bar\epsilon \to 0$,
nên đối số $c-\beta G \to c$; tính liên tục tại $c$ cho
$\Phi(c-\beta G) \to \Phi(c)$. (ii) Tính Lipschitz cục bộ cho
$\|\Phi(c-\beta G) - \Phi(c)\| \le L_\Phi\|\beta G\|_1$; áp dụng lại
Bổ đề~\ref{bd:chuan}.
\end{proof}

\begin{nhanxet}[Khi $c$ cũng thay đổi theo huấn luyện]
Trong thực tế $c$ và $\beta$ được đo lại trên mô hình hiện tại, nên chúng cũng
thay đổi khi $\bar\epsilon \to 0$. Phát biểu đúng khi đó là
$\mathrm{dist}\bigl(\widehat\pi_t, \Phi(c_t)\bigr) \to 0$ dọc theo quá trình
huấn luyện: biên~\eqref{eq:supdo} áp dụng cho từng mô hình, và vế phải tiến
về không miễn $L_{\Phi,t}\beta_t$ bị chặn đều. Mỗi $\Phi(c_t)$ vẫn độc lập
với nhãn, nên kết luận ``đầu ra tiến về một dự đoán độc lập nhãn'' không đổi.
Với ZLG$^+$ và RLU, điều kiện bị chặn đều được bảo đảm vì $\|\pos{c_t}\|_1$
bị chặn dưới (Hệ quả~\ref{hq:apdung}).
\end{nhanxet}

\begin{nhanxet}[Vì sao phải tách hai phần]
Hai phần đòi hỏi \emph{khác hẳn} nhau về mặt tôpô, và đây là điểm dễ nhầm.

Phần (i) chỉ cần tính liên tục \textbf{tại đúng một điểm} $c$: theo định
nghĩa dãy, $x_k\to c$ kéo theo $\Phi(x_k)\to\Phi(c)$, và không đòi gì về
$\Phi$ tại các điểm khác. $\Phi$ có thể gián đoạn tại những điểm gần $c$
\emph{tùy ý} mà kết luận vẫn đúng --- chẳng hạn
$\Phi(x)=\|x-c\|_1(-1)^{\lfloor 1/\|x-c\|_1\rfloor}$ gián đoạn tại vô số
điểm tích tụ về $c$ và không Lipschitz trên bất kì lân cận nào của $c$,
nhưng vẫn cho $\Phi(x_k)\to\Phi(c)=0$.

Phần (ii) thì khác: Lipschitz là một điều kiện \textbf{hai điểm}, nên nó
buộc phải đúng trên cả một \textbf{lân cận} --- ``Lipschitz tại một điểm''
không phải một khái niệm có nghĩa. Đổi lại, nó mua thêm một \emph{tốc độ hội
tụ}, tức một biên dùng được ở $\bar\epsilon$ hữu hạn.

Sự phân biệt này cốt yếu vì Mệnh đề~\ref{md:lipschitz} chứng minh rằng phép
chuẩn hóa mà RLU và ZLG$^+$ dùng \emph{không} Lipschitz \emph{toàn cục};
Bổ đề~\ref{bd:lipcucbo} dưới đây cho thấy nó vẫn Lipschitz \emph{cục bộ}
trên một miền rộng chứa $c$ của hai phương pháp đó
(Nhận xét~\ref{nx:nguongmien} cho ngưỡng tường minh), nên cả (i) và (ii) đều
áp dụng được.
\end{nhanxet}

\begin{bode}[Lipschitz cục bộ của phép cắt-rồi-chuẩn-hóa]
\label{bd:lipcucbo}
Đặt $R(x) = n\,\pos{x}/\|\pos{x}\|_1$. Với mọi $\delta>0$, trên miền
$\{x : \|\pos{x}\|_1 \ge \delta\}$ ánh xạ $R$ là $(2n/\delta)$-Lipschitz theo
chuẩn $L_1$.
\end{bode}

\begin{proof}
Ánh xạ $x \mapsto \pos{x}$ là 1-Lipschitz. Với $v,w \ge 0$ và
$\|v\|_1,\|w\|_1 \ge \delta$,
\[
 \frac{v}{\|v\|_1} - \frac{w}{\|w\|_1}
 = \frac{v-w}{\|v\|_1} + w\,\frac{\|w\|_1-\|v\|_1}{\|v\|_1\,\|w\|_1},
\]
nên, dùng $\bigl|\|w\|_1-\|v\|_1\bigr| \le \|v-w\|_1$ và $\|w\|_1>0$,
\[
 \Bigl\|\tfrac{v}{\|v\|_1} - \tfrac{w}{\|w\|_1}\Bigr\|_1
 \le \frac{\|v-w\|_1}{\|v\|_1} + \frac{\bigl|\|w\|_1-\|v\|_1\bigr|}{\|v\|_1}
 \le \frac{2\|v-w\|_1}{\delta}.
\]
Nhân với $n$ và ghép hai bước.
\end{proof}

\begin{hequa}[ZLG$^+$ và bước đầu của RLU: biên vô điều kiện cho đầu ra]
\label{hq:apdung}
Viết $v = c - \beta\Gamma$ cho đối số thực, với $\Gamma = \tilde G$ cho
ZLG$^+$ và $\Gamma = G$ cho RLU. Đọc đúng thứ tự phép toán trong mã nguồn thì
\begin{equation}
 \textbf{ZLG}^+:\quad \widehat N = \mathrm{round}\bigl(R(\mathrm{int}(v))\bigr),
 \qquad
 \textbf{RLU (lần gọi đầu)}:\quad \widehat N = \mathrm{round}\bigl(R(v)\bigr),
 \label{eq:phantich}
\end{equation}
trong đó $R$ là ánh xạ của Bổ đề~\ref{bd:lipcucbo}. Điểm cần chú ý: ở
ZLG$^+$ phép \texttt{int} nằm \emph{bên trong} phép chuẩn hóa --- nó nhiễu
động \emph{đầu vào} của $R$, không phải đầu ra.
\begin{center}
\footnotesize
\begin{tabular}{@{}llllll@{}}
\toprule
& $c$ & $\beta$ & $\|\pos{c}\|_1$ & đuôi rời rạc & Nguồn \\
\midrule
ZLG$^+$ & $n\bar p$ & $1$ (trên $\tilde G$)
 & $n\sum_j \bar p_j \approx nC$
 & \texttt{int} (trong), \texttt{round} (ngoài) & \texttt{attacks.py:211--217} \\
RLU & $\mathrm{shift}\cdot\bar\rho$ & $1/n$ (trên $G$)
 & $\bar\rho\,\|\mathrm{shift}\|_1 = \bar\rho$
 & \texttt{round} & \texttt{client\_utils.py:128--135} \\
\bottomrule
\end{tabular}
\end{center}
Với ZLG$^+$ giả thiết thêm $w_{\max}$ bị chặn (Mục~\ref{sec:coso}). Với RLU,
phát biểu chỉ áp dụng cho \emph{lần gọi đầu}: lần gọi cuối dùng một $c$ phụ
thuộc cửa sổ (Nhận xét~\ref{nx:gioihan-dl3}). Khi đó:
\begin{enumerate}
\item[(i)] \textup{(Phần thực.)} Vì $\|\pos{\cdot}\|_1$ là $1$-Lipschitz,
 \[
  \|\pos{v}\|_1 \;\ge\; \|\pos{c}\|_1 - \beta\|\Gamma\|_1 \;=:\; \delta ,
 \]
 nên \emph{điều kiện duy nhất} để đối số $v$ nằm trong miền của
 Bổ đề~\ref{bd:lipcucbo} là
 \begin{equation}
  \beta\,\|\Gamma\|_1 \;<\; \|\pos{c}\|_1 .
  \label{eq:dieukienmien}
 \end{equation}
 Khi đó $R$ liên tục và Lipschitz cục bộ với $L_\Phi = 2n/\delta$, và
 $R(v) \to R(c)$ với tốc độ~\eqref{eq:supdo}.
\item[(ii)] \textup{(Biên vô điều kiện trên vector đếm.)} Đặt
 $\gamma_{\mathcal{Q}} = \tfrac12$ cho RLU và
 $\gamma_{\mathcal{Q}} = L_\Phi + \tfrac12$ cho ZLG$^+$. \emph{Không cần
 giả thiết nào thêm,}
 \begin{equation}
  \bigl\|\widehat N - R(c)\bigr\|_1
  \;\le\; \gamma_{\mathcal{Q}}\,C \;+\; L_\Phi\,\beta\,\|\Gamma\|_1
  \;\xrightarrow[\bar\epsilon\to0]{}\; \gamma_{\mathcal{Q}}\,C .
  \label{eq:bienluongtu}
 \end{equation}
\item[(iii)] \textup{(Chuyển sang đại lượng được báo cáo.)} Đặt
 $\mathrm{irec}(M) = n^{-1}\sum_j \min(M_j, N_j)$. Hàm này là
 $1$-Lipschitz theo $\|\cdot\|_1/n$, nên
 \begin{equation}
  \bigl|\mathrm{irec}(\widehat N) - \mathrm{irec}(R(c))\bigr|
  \;\le\; \frac{\gamma_{\mathcal{Q}}\,C + L_\Phi\beta\|\Gamma\|_1}{n}
  \;\xrightarrow[\bar\epsilon\to0]{}\; \frac{\gamma_{\mathcal{Q}}\,C}{n}.
  \label{eq:bienirec}
 \end{equation}
 Với $C = 10$, $n = 400$: RLU cho biên $C/(2n) = 1{,}25\%$; với ZLG$^+$ và
 $\bar p_j \approx 0{,}75$ ta có $\delta = n\sum_j\bar p_j \approx 3000$, nên
 $L_\Phi = 2n/\delta \approx 0{,}27$ và biên là
 $C(L_\Phi+\tfrac12)/n \approx 1{,}92\%$. Tức tỉ lệ khôi phục của hai phương
 pháp bị ghim trong khoảng $2\%$ quanh giá trị của một dự đoán \emph{độc lập
 hoàn toàn với nhãn}.
\end{enumerate}
\end{hequa}

\begin{proof}
(i) là Bổ đề~\ref{bd:lipcucbo} ghép với Định lý~\ref{dl:supdo}.

(ii) Với RLU, bất đẳng thức tam giác qua điểm trung gian $R(v)$:
\[
 \bigl\|\widehat N - R(c)\bigr\|_1
 \le \underbrace{\bigl\|\mathrm{round}(R(v)) - R(v)\bigr\|_1}_{\le\ C/2}
 + \underbrace{\bigl\|R(v) - R(c)\bigr\|_1}_{\le\ L_\Phi \beta \|\Gamma\|_1 \text{ theo (i)}} .
\]
Với ZLG$^+$ cần thêm một bước, và \emph{chính Bổ đề~\ref{bd:lipcucbo} xử lý
nó}: vì $|\mathrm{int}(v_j) - v_j| \le 1$ với mọi $j$, ta có
$\|\mathrm{int}(v)-v\|_1 \le C$, nên
\[
 \bigl\|R(\mathrm{int}(v)) - R(v)\bigr\|_1 \;\le\; L_\Phi\,\|\mathrm{int}(v)-v\|_1
 \;\le\; L_\Phi\,C .
\]
Cộng với sai số của \texttt{round} ($\le C/2$) được
$\gamma_{\mathcal{Q}}C = (L_\Phi+\tfrac12)C$. Cuối cùng áp dụng
Bổ đề~\ref{bd:chuan} cho số hạng $\|\Gamma\|_1$ (với ZLG$^+$, ở dạng có trọng
số của Mục~\ref{sec:coso}).

(iii) Với mỗi $j$, hàm $t \mapsto \min(t, N_j)$ là $1$-Lipschitz, nên
$\bigl|\sum_j \min(a_j,N_j) - \sum_j\min(b_j,N_j)\bigr| \le \sum_j |a_j-b_j|
 = \|a-b\|_1$. Chia cho $n$ và áp dụng (ii).
\end{proof}

\begin{nhanxet}[Miền Lipschitz rộng đến đâu: ngưỡng tường minh]
\label{nx:nguongmien}
Điều kiện~\eqref{eq:dieukienmien} \emph{không} đòi $\beta\Gamma$ phải nhỏ ---
miền của Bổ đề~\ref{bd:lipcucbo} là tập lớn $\{\|\pos{x}\|_1\ge\delta\}$, chứ
không phải một lân cận bé quanh $c$. Thay số với $C=10$, $n=400$:
\begin{center}
\footnotesize
\begin{tabular}{@{}llll@{}}
\toprule
& $\|\pos{c}\|_1$ & $\beta\|\Gamma\|_1$ & Ngưỡng từ~\eqref{eq:dieukienmien} \\
\midrule
ZLG$^+$ & $n\sum_j\bar p_j \approx 3000$ & $\le 2nw_{\max}\bar\epsilon_{\mathrm{tr}} \le 800$
 & thỏa với \emph{mọi} $\bar\epsilon_{\mathrm{tr}}\le1$, trừ khi
   $\bar\epsilon_{\mathrm{aux}}$ gần $1$ \\
RLU & $\bar\rho\,\|\mathrm{shift}\|_1 = \bar\rho \approx 1{,}10$
 & $\le 2\bar\epsilon_{\mathrm{tr}}$
 & $\bar\epsilon_{\mathrm{tr}} < \bar\rho/2 \approx 0{,}55$ \\
\bottomrule
\end{tabular}
\end{center}
Ngưỡng của RLU ứng với một mô hình sai ít nhất $45\%$ số mẫu huấn luyện, tức
rất xa chế độ well-trained mà báo cáo này quan tâm. Vậy
\eqref{eq:bienluongtu}--\eqref{eq:bienirec} \emph{hợp lệ} trên toàn bộ vùng
thực tế.
\end{nhanxet}

\begin{nhanxet}[Hợp lệ không có nghĩa là hữu dụng]
\label{nx:hopleHuudung}
Biên~\eqref{eq:bienirec} đúng trên toàn vùng của
Nhận xét~\ref{nx:nguongmien}, nhưng chỉ \emph{ràng buộc} $\mathrm{irec}$ một
cách có nghĩa khi nó nhỏ hơn hẳn $1$. Thay số ($C=10$, $n=400$,
$\bar\epsilon_{\mathrm{aux}} = 2{,}5\,\bar\epsilon_{\mathrm{tr}}$):
\begin{center}
\footnotesize
\begin{tabular}{@{}lrr@{}}
\toprule
$\bar\epsilon_{\mathrm{tr}}$ & biên~\eqref{eq:bienirec} cho ZLG$^+$ & cho RLU \\
\midrule
$0{,}30$ & $0{,}670$ & $2{,}42$ \\
$0{,}10$ & $0{,}074$ & $0{,}458$ \\
$0{,}05$ & $0{,}041$ & --- \\
$0{,}01$ & $0{,}022$ & $0{,}050$ \\
\bottomrule
\end{tabular}
\end{center}
Vì $\mathrm{irec}\in[0,1]$, một biên $0{,}67$ gần như không nói gì. Biên trở
nên thực sự ràng buộc khi $\bar\epsilon_{\mathrm{tr}}\lesssim 0{,}1$ với
ZLG$^+$ và $\lesssim 0{,}03$ với RLU. Đây là một giới hạn thật và cần nói
kèm: ở vùng $\bar\epsilon_{\mathrm{tr}}$ lớn, định lý vẫn \emph{giải thích cơ
chế} nhưng không cho một biên định lượng hữu ích. Hiện tượng quan sát được
không phụ thuộc vào độ chặt của biên --- Mục~\ref{sec:kiemchung} cho thấy
ZLG$^+$ đã về sát giá trị dự đoán đều ngay từ
$\bar\epsilon_{\mathrm{tr}}\approx0{,}1$.
\end{nhanxet}

\begin{nhanxet}[Vì sao phải viết $\mathcal{Q}$ ra tường minh]
\label{nx:vaitroQ}
Phần (i) \emph{một mình không đủ} để kết luận đầu ra độc lập với nhãn.
Lí do: $R(c-\beta G) = R(c) + e$ với $e$ \emph{phụ thuộc nhãn} và
$\|e\|\to0$; mà ``$\|e\|\to0$'' không kéo theo ``không khôi phục được nhãn
từ $e$'' --- chính VLI$^*$ là phản ví dụ, nó khôi phục \emph{chính xác} từ
tín hiệu cỡ $5\times10^{-12}$ (Mục~\ref{sec:kiemchung}). Suy diễn
``$\|x_k - x\| \to 0 \Rightarrow \|f(x_k)-f(x)\| \to 0$'' \emph{là} định
nghĩa tính liên tục của $f$ tại $x$, nên áp dụng nó cho $\mathcal{Q}$ chính
là ngầm giả định điều cần chứng minh. Vậy $\mathcal{Q}$ là \emph{cái cổng}
quyết định số phận của $e$: xoá nó (đầu ra đóng băng) hay khuếch đại nó
(một chênh lệch nguyên tồn tại mãi). Phần (ii)--(iii) chặn khả năng thứ hai
mà không cần giả thiết nào.
\end{nhanxet}

\begin{nhanxet}[Đóng băng chính xác: một sự tinh chỉnh có điều kiện]
\label{nx:dongbang}
Vì \texttt{int} và \texttt{round} là hằng địa phương, đầu ra còn có thể
\emph{bằng đúng} một vector cố định chứ không chỉ nằm gần nó. Với RLU: nếu
không toạ độ nào của $R(c)$ là số bán nguyên, gọi $\rho$ là khoảng cách tới
số bán nguyên gần nhất; khi $L_\Phi\beta\|G\|_1 < \rho$, mà điều này được
bảo đảm khi $\bar\epsilon < \bar\epsilon_0 = \rho/(2\beta nL_\Phi)$, ta có
$\widehat N = \mathrm{round}(R(c))$ chính xác. Với ZLG$^+$ có \emph{hai}
mép: cần thêm $c_j = n\bar p_j \notin \mathbb Z$ để $\mathrm{int}(v)$ đóng
băng về $\mathrm{int}(c)$, và khi đó giới hạn là
$\mathrm{round}\bigl(R(\mathrm{int}(c))\bigr)$, không phải
$\mathrm{round}(R(c))$. Ba cảnh báo:
\begin{enumerate}
\item \textbf{Ngưỡng không đều theo $c$}: $\bar\epsilon_0 \to 0$ khi
 $R(c)$ tiến tới mép, nên không có một $\bar\epsilon_0$ dùng chung.
\item \textbf{Điều kiện là số học, không phải ``vị trí tổng quát''}: vì tập
 phụ trợ cân bằng theo lớp, $\bar p$ gần đều nên $R(c)_j \approx n/C$; mép
 của \texttt{round} rơi trúng khi $2n/C$ là số nguyên lẻ (với $C=10$ và
 \texttt{batch\_size} chẵn thì không xảy ra). Lập luận ``trừ một tập có độ
 đo không'' \emph{không} áp dụng được, vì $c$ nằm gần đường chéo đều chứ
 không phải một điểm tổng quát.
\item \textbf{Mép của \texttt{int} dễ xảy ra}: chẳng hạn $\bar p_j = 0{,}75$
 và $n = 400$ cho $c_j = 300$ đúng số nguyên, khi đó $\mathrm{int}(v_j)$ ra
 $300$ hoặc $299$ tùy dấu của $\tilde G_j$, tức tùy nhãn. Trong mô phỏng
 chênh lệch đó bị \texttt{round} nuốt mất, nhưng đó là may mắn số học chứ
 không phải điều được chứng minh.
\end{enumerate}
Vì vậy đóng băng chính xác chỉ là một nhận xét, còn
\eqref{eq:bienluongtu}--\eqref{eq:bienirec} mới là phát biểu chính. Mô phỏng
ở Mục~\ref{sec:kiemchung} cho thấy ngay khi $R(c)$ nằm đúng trên mép của
\texttt{round}, độ lệch chỉ là \emph{một} count trên $n$, tức $0{,}25\%$.
\end{nhanxet}

\begin{nhanxet}[Hai phương pháp này \emph{không} cần Mệnh đề~\ref{md:khoangcach}]
Trong Hệ quả~\ref{hq:apdung}, mẫu số $\|\pos{c}\|_1$ bị chặn dưới một cách
\emph{hiển nhiên} ($\bar p \to 1$; $\mathrm{shift}$ là vector xác suất), nên
chỉ cần Bổ đề~\ref{bd:chuan} là đủ để kết luận. Mệnh đề~\ref{md:khoangcach}
về khoảng cách lỗi \emph{không} là trụ cột cho ZLG$^+$ và RLU; nó cần cho
LLG$^+$, nơi $c = \mathrm{offset}$ tự nó là đại lượng $[\epsilon]^1$ và về
nguyên tắc có thể co cùng tốc độ với $G$
(Mệnh đề~\ref{md:llgp}).
\end{nhanxet}

Giới hạn $\Phi(c)$ của hai phương pháp này rất giàu thông tin:

\begin{itemize}
\item Với ZLG$^+$ (dạng~\eqref{eq:zlg}, $c=n\bar p$),
  $\Phi(c) = R(c) = n\bar p/\|\bar p\|_1$. Vì $\bar p_j = 1-\bar\epsilon_{\mathrm{aux},j}
  \to 1$ với \emph{mọi} $j$, vector này xấp xỉ \emph{phân phối đều}, nên ước
  lượng thoái hóa thành việc dự đoán phân phối lớp đều. Dưới phân hoạch
  Dirichlet, độ chồng lấp giữa phân phối thật của client và phân phối đều xấp
  xỉ $0{,}54$ khi $\alpha = 0{,}5$ và $0{,}29$ khi $\alpha = 0{,}1$ với
  $C = 10$; đó chính là các giá trị khôi phục mà phương pháp tiệm cận tới.
  Mục~\ref{sec:kiemchung} xác nhận rằng giới hạn đúng bằng giá trị của dự
  đoán đều. Thất bại ở đây là \emph{tiên đoán được}, không phải ngẫu nhiên.
\item Với RLU (lần gọi đầu của~\eqref{eq:rlu}, $c=\mathrm{shift}\cdot\bar\rho$,
  $\beta=1/n$), $\Phi(c) = n\,\mathrm{shift}/\|\mathrm{shift}\|_1$, cũng là một
  vector độc lập với nhãn.
\end{itemize}

\begin{nhanxet}[Giới hạn của định lý: ba điều nó \emph{không} bao trùm]
\label{nx:gioihan-dl3}
\begin{enumerate}
\item \textbf{LLG$^+$ không có dạng này.} Đầu ra của nó là một vector số
 nguyên sinh bởi vòng lặp $\argminop$, và pha $h_1$ chọn theo \emph{dấu} của
 tín hiệu trước khi trừ offset; nên \emph{không tồn tại} $\Phi$ liên tục để
 áp dụng Định lý~\ref{dl:supdo}. LLG$^+$ được phân tích riêng ở
 Mệnh đề~\ref{md:llgp}: đầu ra cuối cùng chỉ phụ thuộc nhãn qua mẫu dấu, nên
 không \emph{hoàn toàn} độc lập nhãn nhưng mang không quá $C$ bit.
\item \textbf{Thủ tục tìm kiếm hậu nghiệm của RLU không được bao trùm.} Nó
 tính một $\mathrm{shift}$ mới từ ước lượng của lần gọi đầu rồi gọi
 lại~\eqref{eq:rlu}; vì $\mathrm{shift}$ mới phụ thuộc cửa sổ, giả thiết ``$c$
 chỉ phụ thuộc thông tin bên ngoài cửa sổ'' của Định lý~\ref{dl:supdo} không
 còn đúng, nên Hệ quả~\ref{hq:apdung} chỉ phủ \emph{lần gọi đầu}. Mục tiêu
 so khớp của thủ tục chỉ tham chiếu $\mu_{\mathrm{aux}}$ --- một đại lượng
 phụ trợ --- và ước lượng đầu vào của nó đã sụp; tuy vậy nó \emph{có} tiêu
 thụ thêm một đại lượng từ cửa sổ:
 $\|O\|^2$ với $O_d = \overline{(\Delta W_{\cdot d}/\Delta b)}$, một
 \emph{tỉ số} gradient nên bất biến thang độ. Đó là chỗ duy nhất RLU về
 nguyên tắc còn giữ được thông tin của cửa sổ; vì $\|O\|^2$ chỉ là một vô
 hướng điều tiết độ dài bước, chúng tôi không cho rằng nó phục hồi được số
 lượng theo lớp, nhưng một chứng minh chặt cần phân tích vòng lặp và chúng
 tôi không đưa ra ở đây.
\item \textbf{iLRG không thuộc lớp này} vì nó không dùng thống kê phụ trợ
 nào; xem Mục~\ref{sec:ilrg}.
\end{enumerate}
\end{nhanxet}

\subsection{Vì sao độ lớn ngoại lai không co lại cùng tín hiệu}

\begin{menhde}[Khoảng cách lỗi trên dữ liệu chưa thấy]
\label{md:khoangcach}
Nếu một mẫu bị phân loại sai theo quy tắc $\argmaxop$, thì
$p_{i,y_i} \le 1/2$ và do đó $\epsilon_i \ge 1/2$. Hệ quả là, nếu tỉ lệ lỗi
trên tập phụ trợ là $\tau_{\mathrm{err}}$, thì
\[
 \bar\epsilon_{\mathrm{aux}} \;\ge\; \tau_{\mathrm{err}}/2 .
\]
Do đó
\begin{equation}
 \frac{\bar\epsilon_{\mathrm{aux}}}{\bar\epsilon_{\mathrm{tr}}}
 \;\ge\; \frac{\tau_{\mathrm{err}}}{2\,\bar\epsilon_{\mathrm{tr}}},
 \label{eq:tiso}
\end{equation}
Vì vậy một \emph{điều kiện đủ} để tỉ số bên trái phân kỳ là
$\tau_{\mathrm{err}}/\bar\epsilon_{\mathrm{tr}} \to \infty$, tức lỗi trên dữ
liệu chưa thấy giảm \emph{chậm hơn} lỗi huấn luyện. Điều kiện này \emph{không
cần thiết}: tỉ số vẫn có thể phân kỳ khi $\tau_{\mathrm{err}} = 0$, nếu các
mẫu phụ trợ đúng nhưng có margin thấp.
\end{menhde}

\begin{proof}
Trên một mẫu bị phân loại sai, tồn tại $k \ne y_i$ với
$p_{i,k} \ge p_{i,y_i}$; vì $p_{i,k} + p_{i,y_i} \le 1$ nên
$p_{i,y_i} \le 1/2$. Lấy trung bình biên dưới thu được cho $1 - p_{i,y_i}$
trên toàn tập phụ trợ, bỏ đi phần đóng góp không âm của các mẫu phân loại
đúng.
\end{proof}

\begin{nhanxet}[Điều kiện là \emph{tương đối}, và nó là một giả thiết]
\label{nx:tiso}
Điều kiện đủ ở trên \emph{không} đòi $\tau_{\mathrm{err}}$ lớn, cũng
không đòi nó bị chặn dưới bởi một hằng số tuyệt đối. Nó chỉ đòi
$\tau_{\mathrm{err}}$ giảm chậm hơn $\bar\epsilon_{\mathrm{tr}}$ --- và đó
chính là định nghĩa của một khoảng cách tổng quát hóa. Điều kiện này rất dễ
thỏa: ngay cả với $\tau_{\mathrm{err}} = 0{,}003$ (mức MNIST rất tốt), nếu
$\bar\epsilon_{\mathrm{tr}} = 10^{-4}$ thì~\eqref{eq:tiso} vẫn đảm bảo tỉ số
$\ge 15$; với $\bar\epsilon_{\mathrm{tr}} = 10^{-6}$ thì $\ge 1500$. Lí do là
lỗi huấn luyện có thể bị đẩy về \emph{đúng} không bằng ghi nhớ, còn lỗi kiểm
tra thì bình nguyên.

Tuy vậy đây vẫn là một \emph{giả thiết về mô hình và dữ liệu}, không phải một
kết quả. Nếu $\tau_{\mathrm{err}} \asymp \bar\epsilon_{\mathrm{tr}}$, chẳng
hạn $\tau_{\mathrm{err}} = c\,\bar\epsilon_{\mathrm{tr}}$, thì~\eqref{eq:tiso}
chỉ cho tỉ số $\ge c/2$: biên \emph{không còn đảm bảo} phân kỳ, dù tỉ số thật
vẫn có thể phân kỳ nhờ các mẫu đúng có margin thấp. Điều này chỉ ảnh hưởng
tới LLG$^+$ --- ZLG$^+$ và RLU không cần mệnh đề này --- và cho một tiên đoán
khả phủ chứng, nêu ở Mục~\ref{sec:gioihan}.

Ngoài ra biên là \emph{bảo thủ}: nó chỉ dùng các mẫu bị phân loại sai và bỏ
qua hoàn toàn đóng góp của các mẫu đúng nhưng margin thấp. Vậy
$\bar\epsilon_{\mathrm{aux}}$ thực tế \emph{lớn hơn} mức biên nói, điều này
củng cố chứ không làm yếu kết luận.
\end{nhanxet}

Mệnh đề~\ref{md:khoangcach} là nửa thứ hai của cơ chế cho \emph{LLG$^+$}:
Bổ đề~\ref{bd:chuan} nói $\|G\| \to 0$, còn ở đó $c = \mathrm{offset}$ tự nó
là đại lượng $[\epsilon]^1$ nên về nguyên tắc có thể co cùng tốc độ ---
\eqref{eq:tiso} là thứ bịt lối đó. Với ZLG$^+$ và RLU thì không cần, vì
$\|\pos{c}\|_1$ đã bị chặn dưới hiển nhiên (Hệ quả~\ref{hq:apdung}).

\subsection{Thất bại là \emph{tương tác}, không phải riêng việc tín hiệu co nhỏ}
\label{sec:tuongtac}

Sẽ sai nếu kết luận ``chính việc tín hiệu co nhỏ --- hay khoảng cách tổng
quát hóa --- phá vỡ các tấn công này''. Việc co nhỏ \emph{không} phá hủy
thông tin: nó chỉ làm thành phần mang nhãn trở nên nhỏ so với $c$. Mệnh đề
sau cho thấy thông tin vẫn khôi phục được \emph{từ đại lượng thực $v$ mà
ZLG$^+$ tính ra trước khi lượng tử hóa}.

\begin{menhde}[Thông tin nhãn còn nguyên trong $R(v)$]
\label{md:conthongtin}
Giả sử $c>0$ theo từng thành phần, và đặt $H=\{x:\one^Tx=0\}$. Trên $H$,
nhân Jacobi của $R$ tại $c$ tác động như một vô hướng:
\[
 DR(c)\,x = \frac{n}{\|c\|_1}\,x \qquad (x \in H).
\]
Vì tín hiệu $\Gamma$ (là $G$ hoặc $\tilde G$) nằm trong $H$
theo~\eqref{eq:tongkhong}, ta có
\[
 e := R(c-\beta\Gamma) - R(c) = -\frac{\beta n}{\|c\|_1}\,\Gamma + o(\|\Gamma\|):
\]
ở bậc nhất, $e$ là một bội \emph{âm} của chính tín hiệu. Do đó $\Gamma$ khôi
phục được từ $e$ tới một hệ số dương, và mọi ước lượng bất biến thang độ ---
chẳng hạn Định lý~\ref{dl:dangdong}, dưới giả thiết của nó --- cho cùng kết
quả khi áp dụng lên $-e$ hay lên $\Gamma$.
\end{menhde}

\begin{proof}
Với $x>0$ ta có $R(x) = n x/\one^Tx$, nên
$DR(c) = n\bigl(I/\|c\|_1 - c\,\one^T/\|c\|_1^2\bigr)$. Với $x\in H$ thì
$\one^Tx = 0$, nên số hạng thứ hai triệt tiêu. Khai triển Taylor bậc nhất cho
phần còn lại.
\end{proof}

\begin{nhanxet}[Kẻ tấn công \emph{biết} $c$, nên $e$ là đại lượng tính được]
\label{nx:doclaiti}
$c$ được chính kẻ tấn công đo trên tập phụ trợ của họ, nên $e = R(v)-R(c)$ là
tính được. Mục~\ref{sec:kiemchung} cho thấy khi đọc $-e$ \emph{theo tỉ số}
(không thang độ), tỉ lệ khôi phục là $1{,}000$ --- \emph{chính xác tuyệt
đối} --- ngay tại $\|e\|_\infty \approx 2{,}4\times10^{-7}$, trong khi
ZLG$^+$ trên cùng dữ liệu đó chỉ cho đúng giá trị của dự đoán đều. Nói cách
khác, thông tin chưa bao giờ rời khỏi cập nhật; ZLG$^+$ đánh mất nó vì cộng
một nền lớn rồi lượng tử hóa.
\end{nhanxet}

Vậy phát biểu đúng là một \textbf{hội}:
\[
 \text{thất bại}
 \;=\;
 \underbrace{\text{lệch thang độ } \|c\| \gg \beta\|\Gamma\|}_{\text{xuất hiện khi } \bar\epsilon_{\mathrm{tr}}\to0}
 \;\wedge\;
 \underbrace{\text{thiết kế phụ thuộc thang độ}}_{\text{so độ lớn tuyệt đối, rồi lượng tử hóa}} .
\]
Lệch thang độ có \emph{hai nguồn khác nhau}. Với ZLG$^+$ và RLU, $c$ có cỡ
một vector xác suất ($\bar p\approx1$; $\mathrm{shift}$ có tổng $1$), nên chỉ
cần $\bar\epsilon_{\mathrm{tr}}\to0$; khoảng cách tổng quát hóa không đóng
vai trò gì. Với LLG$^+$, $c=\mathrm{offset}$ tự nó có cỡ
$\bar\epsilon_{\mathrm{aux}}$, nên lệch thang độ đòi
$\bar\epsilon_{\mathrm{aux}}\gg\bar\epsilon_{\mathrm{tr}}$ --- đây mới là chỗ
khoảng cách tổng quát hóa (Mệnh đề~\ref{md:khoangcach}) đi vào. Không vế nào
của hội một mình đủ: nếu $\|c\|\sim\beta\|\Gamma\|$ thì các thiết kế đó vẫn
chạy; còn nếu ước lượng bất biến thang độ thì lệch thang độ không ảnh hưởng
--- đó chính là VLI$^*$ (Mệnh đề~\ref{md:batbien}). Hệ quả thực tiễn: cách
sửa \emph{không phải} ước lượng phụ trợ tốt hơn mà là \textbf{bất biến thang
độ}. Mệnh đề~\ref{md:conthongtin} cũng cho một đối chứng kiểm được: đọc cùng
tín hiệu đó theo tỉ số thì \emph{không} suy giảm.

\subsection{LLG\texorpdfstring{$^+$}{+}: chỉ còn phụ thuộc nhãn qua mẫu dấu}
\label{sec:llgp}

Pha $h_1$ của LLG$^+$ chọn $\{j : u_j < 0\}$ trên tín hiệu $u$ \emph{trước
khi} trừ offset. Tập này chỉ phụ thuộc \emph{dấu} của $u$, mà dấu không đổi
khi nhân $u$ với một hệ số dương. Vì vậy $h_1$ không tiêu biến khi
$\bar\epsilon\to0$, và Định lý~\ref{dl:supdo} không mô tả được LLG$^+$. Nhưng
$h_1$ cũng là \emph{tất cả} những gì còn lại.

\begin{menhde}[LLG$^+$: đầu ra là hàm của mẫu dấu]
\label{md:llgp}
Xét LLG$^+$ như cài đặt, với $\mathrm{offset}$ và $\mathrm{impact}<0$ đo trên
dữ liệu phụ trợ (coi là cố định). Với mỗi tập $h \subseteq \{1,\dots,C\}$, gọi
$F(h)$ là đầu ra khi chạy pha $h_1 = h$ rồi vòng tham lam $B-|h|$ bước xuất
phát từ $g^0(h) = -\mathrm{offset} + |\mathrm{impact}|\,\one_h$, và gọi
$\Delta(h)$ là hiệu nhỏ nhất giữa hai toạ độ được so sánh tại mọi bước của
lần chạy đó. Giả thiết lần chạy đó \emph{không có hòa}, tức $\Delta(h)>0$.
\begin{enumerate}
\item[(i)] \textup{($h_1$ bất biến thang độ.)} $h_1(u) = h_1(cu)$ với mọi
 $c>0$. Với tín hiệu bias $G$ thì theo Bổ đề~\ref{bd:dongnhat},
 $G_j<0 \iff a_j > E_j$; với tín hiệu có trọng số $\tilde G$ là phiên bản có
 trọng số của điều kiện này. Riêng toạ độ có $u_j = 0$ đúng là một mép: dấu
 của nó không xác định.
\item[(ii)] \textup{(Phần còn lại đóng băng.)} Nếu
 $2\|u\|_\infty < \Delta(h_1(u))$ thì đầu ra \emph{bằng chính xác}
 $F(h_1(u))$.
\end{enumerate}
Do đó khi $\bar\epsilon_{\mathrm{tr}}$ đủ nhỏ, đầu ra chỉ phụ thuộc nhãn qua
$h_1$, nên nhận không quá $2^C$ giá trị và mang không quá $C$ bit thông tin
về nhãn, so với $\log_2\binom{n+C-1}{C-1}$ bit cần để xác định $N$ (khoảng
$59{,}5$ bit khi $n=400$, $C=10$).
\end{menhde}

\begin{proof}
(i) $\mathrm{sign}(cu_j) = \mathrm{sign}(u_j)$ với $c>0$. (ii) Khi chạy với
tín hiệu $u$, vector bước vào vòng tham lam là $g^0(h_1) + u$, và chừng nào các
lựa chọn còn trùng với lần chạy tham chiếu thì mọi bước cộng
$|\mathrm{impact}|$ vào cùng toạ độ. Vậy mỗi toạ độ lệch khỏi lần chạy tham
chiếu đúng một lượng $u_k$, nên hiệu giữa hai toạ độ bất kì lệch không quá
$2\|u\|_\infty < \Delta(h_1)$; không cặp nào đổi thứ tự, và bằng quy nạp theo
bước, dãy $\argminop$ trùng với lần chạy tham chiếu. (Mã nguồn dùng hai giá
trị impact phụ trợ hơi khác nhau cho hai pha; điều đó không ảnh hưởng tới lập
luận.)
\end{proof}

Mệnh đề~\ref{md:khoangcach} đi vào đúng chỗ này: $\Delta$ có cỡ của
$\mathrm{offset}$ và $|\mathrm{impact}|$, tức cỡ $\bar\epsilon_{\mathrm{aux}}$,
còn $\|u\|_\infty$ có cỡ $\bar\epsilon_{\mathrm{tr}}$, nên điều kiện ở (ii)
đòi $\bar\epsilon_{\mathrm{tr}} \ll \bar\epsilon_{\mathrm{aux}}$. Trong mô
phỏng ở Mục~\ref{sec:kiemchung}, đầu ra khớp $F(h_1)$ trong $500/500$ trường
hợp, và trên $3000$ vector nhãn ngẫu nhiên chỉ có $4$ đầu ra khác nhau ---
tức LLG$^+$ thực tế gần như sụp hoàn toàn, với irec còn thấp hơn dự đoán đều.
Điều này cũng giải thích vì sao LLG thuần bền vững: nó dùng cùng pha $h_1$,
\emph{và} $\mathrm{impact}$ của nó được tính từ chính tín hiệu nên vòng tham
lam của nó cũng bất biến thang độ.

\subsection{Hai khiếm khuyết riêng của RLU}

\begin{menhde}[Khoảng cách Jensen khi lấy trung bình logit]
\label{md:jensen}
Với một mẫu có margin $M$, ta có $\epsilon = (1+e^M)^{-1}$, và
$\epsilon \approx e^{-M}$ khi $M$ lớn và dương. Nếu margin trong một lớp xấp
xỉ $\mathcal N(\mu,\sigma^2)$ và xấp xỉ margin lớn còn hợp lệ trên phần lớn
khối lượng của nó, thì
\[
 \E\bigl[e^{-M}\bigr] = e^{-\mu + \sigma^2/2},
 \qquad e^{-\E[M]} = e^{-\mu},
\]
nên việc thay bằng margin trung bình \emph{đánh giá thấp} đại lượng thay thế
này đi một hệ số $e^{\sigma^2/2}$. Bất đẳng thức Jensen cho
$\E[e^{-M}] \ge e^{-\E[M]}$ một cách vô điều kiện.
\end{menhde}

Phương trình~\eqref{eq:rlu} đặt
$\mathrm{shift} = \softmaxop(\mu_{\mathrm{aux}})$, tức softmax của
\emph{logit trung bình} --- chính là phép thay thế mà
Mệnh đề~\ref{md:jensen} mô tả. Độ chệch là hệ thống và theo chiều
\emph{đánh giá thấp} khối lượng lỗi. Chi tiết làm nặng thêm vấn đề: huấn
luyện làm tăng \emph{cả} $\mu$ \emph{và} $\sigma$, nên khoảng cách
$e^{\sigma^2/2}$ lớn nhất đúng vào lúc tín hiệu yếu nhất; hai hiệu ứng cộng
hưởng.

\begin{menhde}[$R$ không Lipschitz toàn cục; phép chiếu thì có]
\label{md:lipschitz}
Đặt $R(y) = n\pos{y}/\|\pos{y}\|_1$ (phép cắt rồi chuẩn hóa, dùng
trong~\eqref{eq:rlu} và sau~\eqref{eq:zlg}), và gọi $\Pi_{\simplex}$ là phép
chiếu Euclid lên $\simplex$. Khi đó $R$ không có biên Lip\-schitz đều nào,
còn $\Pi_{\simplex}$ thì có:
\begin{gather*}
 \sup_{a \ne b} \frac{\|R(a)-R(b)\|_1}{\|a-b\|_1} = \infty, \\[0.2em]
 \|\Pi_{\simplex}(a) - \Pi_{\simplex}(b)\|_2 \le \|a-b\|_2
 \quad \text{với mọi } a,b,
\end{gather*}
và do tương đương chuẩn, $\|\Pi_{\simplex}(a)-\Pi_{\simplex}(b)\|_1 \le
\sqrt{C}\,\|a-b\|_1$. Điểm cốt yếu là $R$ \emph{không} có biên đều nào, còn
$\Pi_{\simplex}$ thì có, trong cả hai chuẩn.
\end{menhde}

\begin{proof}
Phép chiếu lên một tập lồi đóng khác rỗng là không giãn trong chuẩn
\emph{Euclid} --- chuẩn định nghĩa phép chiếu --- cho khẳng định thứ hai;
biên $L_1$ suy ra từ $\|x\|_1 \le \sqrt C\|x\|_2$ và $\|x\|_2 \le \|x\|_1$.
Với khẳng định thứ nhất, lấy
$a = (\delta, \delta/2, -1, \ldots, -1)$ và nhiễu theo hướng $e_1$. Tại $a$ ta
có $\|\pos{a}\|_1 = 1{,}5\,\delta$ và $R(a)/n = (2/3, 1/3, 0, \ldots)$, nên
đạo hàm theo hướng là
\[
 \lim_{\eta\to0}\frac{\|R(a+\eta e_1)-R(a)\|_1}{\eta}
 = \frac{n}{1{,}5\,\delta}\,\bigl\|e_1 - R(a)/n\bigr\|_1
 = \frac{4n}{9\delta} \xrightarrow[\delta\to0]{} \infty .
\]
Hướng nhiễu phải không song song với $\pos{a}$; nếu song song thì $R$ không
đổi.
\end{proof}

Hệ quả thực tiễn là bước chuẩn hóa cuối của RLU và ZLG$^+$ có thể
\emph{khuếch đại} nhiễu gradient không giới hạn, còn bước tương ứng của
VLI$^*$ (Mục~\ref{sec:vli}) thì không thể.

\begin{nhanxet}[Không có ``hệ số khuếch đại'' nào để báo cáo]
\label{nx:khongtrichso}
Vì supremum bằng vô cùng, \emph{mọi} con số cụ thể lấy từ một lần lấy mẫu đều
vô nghĩa: nó chỉ phản ánh khoảng lấy mẫu chứ không phải một tính chất của $R$.
Trong hai lần quét $2\times10^5$ cặp điểm với hai khoảng khác nhau, chúng tôi
thu được $4{,}4\times10^4$ và $3{,}4\times10^5$. Cách báo cáo đúng là công
thức đạo hàm theo hướng $4n/(9\delta)$: nó khớp giá trị đo tới ba chữ số có
nghĩa trên bốn bậc độ lớn của $\delta$ (Mục~\ref{sec:kiemchung}), và cho thấy
rõ sự phân kỳ.
\end{nhanxet}

\begin{nhanxet}[Toàn cục so với cục bộ: không có mâu thuẫn]
Mệnh đề~\ref{md:lipschitz} là một phát biểu \emph{toàn cục}; nó không xung
đột với Bổ đề~\ref{bd:lipcucbo}, vốn cho tính Lipschitz \emph{cục bộ} với
hằng số $2n/\delta$ trên miền $\|\pos{x}\|_1 \ge \delta$. Hai kết quả đóng
hai vai khác nhau: Bổ đề~\ref{bd:lipcucbo} cho phép \emph{áp dụng}
Định lý~\ref{dl:supdo} cho ZLG$^+$ và RLU (Hệ quả~\ref{hq:apdung}); còn
Mệnh đề~\ref{md:lipschitz} nói rằng hằng số cục bộ ấy \emph{nổ} khi
$\|\pos{x}\|_1 \to 0$, nên không có biên đều nào trên toàn miền --- đó là
rủi ro khuếch đại nhiễu mà VLI$^*$ tránh được.
\end{nhanxet}

\subsection{iLRG: lỗi điều kiện số, không phải lỗi hiệu chuẩn}
\label{sec:ilrg}

iLRG không dùng thống kê phụ trợ nào, nên Định lý~\ref{dl:supdo} không áp
dụng. Khó khăn của nó khác. Theo Bổ đề~\ref{bd:dongnhat},
$G_j = -a_j + E_j$, nên khi $\bar\epsilon \to 0$ thì \emph{cả} $a_j$ và $E_j$
đều tiêu biến, và gradient bias dùng làm \emph{mẫu số} trong
$\hat e_j = \nabla_{W_j}/\nabla_{b_j}$ tiến về không. Tỉ số vẫn hữu hạn trong
số học chính xác, nhưng sai số tương đối của nó bị khuếch đại không giới hạn;
cài đặt tham chiếu buộc phải chặn số mũ của tỉ số và thay bằng một giá trị
định sẵn khi mẫu số đúng bằng không. Hệ tuyến tính sau đó cũng ngày càng suy
biến, vì $p_i \to e_{y_i}$ làm mọi vector xác suất khôi phục collapse về các
vector one-hot.

\subsection{Tóm tắt chẩn đoán}

\begin{center}
\footnotesize
\begin{tabular}{@{}lllp{6.0cm}@{}}
\toprule
Phương pháp & Căn cứ & Sụp đổ? & Giới hạn hoặc triệu chứng \\
\midrule
iLRG & BĐ~\ref{bd:dongnhat} (điều kiện số) & --- &
 chia cho gradient bias tiêu biến; giả nghịch đảo gần kỳ dị \\
LLG & không áp dụng & \textbf{không} &
 $h_1$ theo dấu \emph{và} impact tính từ $G$: bất biến thang độ. Nhưng giả
 định $r_j\equiv1$, $Q$ đều, và nhân đáp án một batch với $m$ \\
LLG$^+$ & MĐ~\ref{md:llgp} (+ MĐ~\ref{md:khoangcach}) & \textbf{gần toàn phần} &
 đầu ra chỉ còn là hàm của mẫu dấu $h_1$ ($\le C$ bit); mô phỏng: $4$ đầu ra
 khác nhau, irec dưới mức dự đoán đều \\
ZLG$^+$ & ĐL~\ref{dl:supdo} + HQ~\ref{hq:apdung} & \textbf{toàn phần} &
 $\widehat N \to n\bar p/\|\bar p\|_1 \approx$ phân phối đều (cần $w_{\max}$ bị chặn) \\
RLU & ĐL~\ref{dl:supdo} + HQ~\ref{hq:apdung}, MĐ~\ref{md:jensen},~\ref{md:lipschitz} & \textbf{toàn phần} (lần gọi đầu) &
 $\widehat N \to n\,\mathrm{shift}/\|\mathrm{shift}\|_1$; độ chệch Jensen
 $e^{\sigma^2/2}$; chuẩn hóa khuếch đại nhiễu \\
\bottomrule
\end{tabular}
\end{center}

% ============================================================================
\section{VLI và VLI$^*$}
\label{sec:vli}

\subsection{Nguyên lý thiết kế}

Bổ đề~\ref{bd:dongnhat} viết $a_j = N_j r_j \bar\epsilon$, nên thông tin nhãn
nằm trong cấu trúc \emph{tương đối} của $G$, còn $\bar\epsilon$ --- đại lượng
tiêu biến khi huấn luyện --- chỉ đi vào như một hệ số chung. Điều này gợi ra
việc đòi hỏi ở ước lượng đúng cái tính chất mà các phương pháp cũ thiếu.

\begin{menhde}[Bất biến thang độ]
\label{md:batbien}
Gọi $\Psi$ là ánh xạ định nghĩa trong Định lý~\ref{dl:dangdong}, và
$\Psi_{\mathrm{tq}}$ là ánh xạ của Định lý~\ref{dl:tongquat}. Với mọi $c>0$,
\[
 \Psi(cG) = \Psi(G), \qquad
 \Psi_{\mathrm{tq}}(cG, r, Q) = \Psi_{\mathrm{tq}}(G,r,Q) .
\]
Do đó không ước lượng nào trong hai có thể bộc lộ sự sụp đổ của
Định lý~\ref{dl:supdo}: không đại lượng có độ lớn tuyệt đối nào đi vào hai
ánh xạ này.
\end{menhde}

\begin{proof}
Trong Định lý~\ref{dl:dangdong}, tín hiệu chỉ xuất hiện qua tỉ số
$G_j/G_{\max}$, và $G_{\max}(cG) = c\,G_{\max}(G)$ với $c>0$. Trong
Định lý~\ref{dl:tongquat}, nghiệm $a$ của~\eqref{eq:hept} là thuần nhất dương
theo $G$, và đầu ra là tỉ số $n(a_j/r_j)/\sum_k (a_k/r_k)$.
\end{proof}

\begin{nhanxet}[Bản heuristic chỉ bất biến một phần]
Thủ tục VLI ban đầu dựng một cơ sở tổng-bằng-không với đường chéo
$d = n^{-1}\sum_{j:G_j<0}G_j$ và ngoài đường chéo $-d/(C-1)$, nhân $d$ với
một ``dynamic boost'' $f_1 f_2 f_3$, rồi lần lượt bóc các cột cơ sở khỏi một
phần dư khởi tạo bằng $G$. Hệ số $f_1$ là hàm của các hình chiếu đã chuẩn hóa
nên bất biến thang độ, nhưng $f_2 = (w_2\|G\|_2+1)^{-1}$ và
$f_3 = w_3^{-100|d|}$ \emph{phụ thuộc độ lớn tuyệt đối} của $G$. Vậy
heuristic vi phạm Mệnh đề~\ref{md:batbien}, và năm hằng số của nó phải dò lại
cho từng tập dữ liệu, từng mô hình và từng mức độ chính xác.
Định lý~\ref{dl:dangdong} loại bỏ chúng.
\end{nhanxet}

\subsection{Giả thiết đối xứng, và vì sao nó bị cưỡng bức}

\begin{giathiet}[Hình học lỗi đối xứng]
\label{gt:doixung}
Mọi lớp có mặt đều có lỗi trung bình bằng nhau, $r_j = 1$, và mỗi lớp phân bổ
khối lượng lỗi của nó đều trên các lớp còn lại,
$Q_{jk} = 1/(C-1)$ với $k \ne j$. Có ít nhất một lớp vắng mặt khỏi cửa sổ, và
$\bar\epsilon > 0$.
\end{giathiet}

Giả thiết~\ref{gt:doixung} là \emph{sai} với mọi mô hình thực. Mệnh đề sau
giải thích vì sao nó vẫn là lựa chọn đúng cho một ước lượng bị đòi hỏi không
dùng thông tin riêng về lớp.

\begin{menhde}[Tính duy nhất dưới tính hoán vị nhãn]
\label{md:hoanvi}
Ma trận đều $Q_{jk} = 1/(C-1)$ $(k\ne j)$, $Q_{jj}=0$, là ma trận
\emph{duy nhất} vừa row-stochastic, vừa có đường chéo bằng không, vừa thỏa
$P Q P^T = Q$ với mọi ma trận hoán vị $P$ cỡ $C\times C$.
\end{menhde}

\begin{proof}
Bất biến hoán vị nghĩa là $Q_{\sigma(j)\sigma(k)} = Q_{jk}$ với mọi hoán vị
$\sigma$. Cho hai vị trí ngoài đường chéo bất kì $(j,k)$ và $(j',k')$ với
$j\ne k$, $j'\ne k'$, tồn tại một hoán vị đưa vị trí này sang vị trí kia, nên
mọi phần tử ngoài đường chéo bằng nhau, giả sử bằng $q$. Tính row-stochastic
với $C-1$ phần tử ngoài đường chéo mỗi hàng buộc $q = 1/(C-1)$. Ngược lại, ma
trận đều rõ ràng bất biến.
\end{proof}

Một kẻ tấn công không biết lớp nào dễ lẫn với lớp nào thì không có căn cứ nào
để phá vỡ tính đối xứng giữa các nhãn lớp; mọi $Q$ không đều đều ngầm khẳng
định một tri thức riêng về lớp. Mệnh đề~\ref{md:hoanvi} nói rằng khi đó lựa
chọn đều không phải là một sự đơn giản hóa mà là lựa chọn \emph{duy nhất chấp
nhận được} --- cùng kiểu lập luận suy ra phân phối tiên nghiệm đều từ tính
hoán vị được. Việc đếm bậc tự do cũng cho kết luận tương tự: hệ~\eqref{eq:hept}
cho $C$ phương trình, trong khi các ẩn là $a$ cùng với $C(C-2)$ tham số tự do
của $Q$; để định danh được mà không có thông tin phụ, phải triệt tiêu toàn bộ
nhóm sau, và ma trận đều là lựa chọn duy nhất không có tham số tự do, tương
đẳng với phép đổi tên lớp, và vẫn giữ hệ giải được ở dạng đóng.

\subsection{Dạng đóng}

\begin{dinhly}[Dạng đóng VLI$^*$]
\label{dl:dangdong}
Dưới Giả thiết~\ref{gt:doixung},
\begin{equation}
 G_j = \frac{n\bar\epsilon - C a_j}{C-1},
 \qquad
 G_{\max} = \frac{n\bar\epsilon}{C-1},
 \qquad
 \boxed{\;N_j = \frac{n}{C}\Bigl(1 - \frac{G_j}{G_{\max}}\Bigr)\;}
 \label{eq:dangdong}
\end{equation}
và biểu thức cuối bất biến dưới $G \mapsto cG$ với $c>0$.
\end{dinhly}

\begin{proof}
Dưới nhầm lẫn đều, $(Q^Ta)_j = \bigl(\sum_k a_k - a_j\bigr)/(C-1)$, nên
theo~\eqref{eq:hept}
\[
 G_j = -a_j + \frac{\sum_k a_k - a_j}{C-1}
     = \frac{\sum_k a_k - C a_j}{C-1} ,
\]
và $\sum_k a_k = n\bar\epsilon$ cho đẳng thức thứ nhất. Vì $a_j \ge 0$ và
$G_j$ là hàm affine giảm nghiêm ngặt theo $a_j$ với \emph{cùng} độ dốc
$-C/(C-1)$ cho mọi $j$, cực đại của $G$ đạt được đúng tại nơi $a_j = 0$, tức
tại một lớp vắng mặt, cho $G_{\max} = n\bar\epsilon/(C-1)$;
Giả thiết~\ref{gt:doixung} bảo đảm lớp như vậy tồn tại. Chia hai vế,
\[
 \frac{G_j}{G_{\max}} = 1 - \frac{C a_j}{n\bar\epsilon}
 \quad\Longrightarrow\quad
 a_j = \frac{n\bar\epsilon}{C}\Bigl(1 - \frac{G_j}{G_{\max}}\Bigr),
\]
và $r_j = 1$ cho $a_j = N_j\bar\epsilon$, nên $\bar\epsilon$ triệt tiêu và ta
được~\eqref{eq:dangdong}. Một hệ số dương chung triệt tiêu trong tỉ số
$G_j/G_{\max}$.
\end{proof}

Việc $\bar\epsilon$ triệt tiêu ở bước cuối chính là toàn bộ ý nghĩa. Vai trò
của $G_{\max}$ là làm một \emph{nhiệt kế nằm trong cửa sổ} để đo
$\bar\epsilon$: một lớp vắng mặt có $a_j = 0$, nên thành phần của nó đọc
thang độ lỗi trực tiếp từ chính cập nhật bị tấn công. Ở chỗ RLU, ZLG$^+$ và
LLG$^+$ so tín hiệu với một đại lượng đo từ dữ liệu phụ trợ --- và bị
Định lý~\ref{dl:supdo} hoặc Mệnh đề~\ref{md:llgp} trừng phạt khi tín hiệu co
nhỏ --- VLI$^*$ đọc thang độ từ bên trong cửa sổ.

\begin{menhde}[Tính chất cấu trúc của~\eqref{eq:dangdong}]
\label{md:tinhchat}
Đặt $\mathrm{raw}_j = \tfrac{n}{C}(1 - G_j/G_{\max})$ với $G_{\max}>0$. Khi
đó
\begin{enumerate}
\item $\mathrm{raw}_j \ge 0$ với mọi $j$, đẳng thức tại thành phần đạt cực
  đại; \emph{không cần} phép cắt.
\item $\sum_j \mathrm{raw}_j = n$, theo~\eqref{eq:tongkhong};
  \emph{không cần} phép chuẩn hóa.
\item Dưới Giả thiết~\ref{gt:doixung},
  $\mathrm{raw}_j - \mathrm{raw}_k = N_j - N_k$ một cách chính xác.
\end{enumerate}
\end{menhde}

\begin{proof}
(i) $G_j \le G_{\max}$ và $G_{\max}>0$ cho $G_j/G_{\max} \le 1$.
(ii) $\sum_j \mathrm{raw}_j = \tfrac{n}{C}\bigl(C - G_{\max}^{-1}\sum_j G_j\bigr) = n$.
(iii) Từ Định lý~\ref{dl:dangdong},
$G_k - G_j = C\bar\epsilon(N_j-N_k)/(C-1)$ và
$G_{\max} = n\bar\epsilon/(C-1)$, nên
$\mathrm{raw}_j - \mathrm{raw}_k = \tfrac{n}{C}\cdot\tfrac{G_k-G_j}{G_{\max}}
 = \tfrac{n}{C}\cdot\tfrac{C(N_j-N_k)}{n} = N_j - N_k$.
\end{proof}

Hai ý (i) và (ii) loại bỏ đúng hai bước mà Mệnh đề~\ref{md:lipschitz} chỉ ra
là bệnh lý trong RLU và ZLG$^+$. Ý (iii) xác định cách sửa đúng cho một ước
lượng bị nhiễu: thông tin nhãn nằm trong \emph{hiệu} của $\mathrm{raw}$, nên
phép sửa phải là phép \emph{cộng}. Phép chiếu Euclid lên $\simplex$ có đúng
dạng đó,
\[
 \Pi_{\simplex}(y)_j = \max(y_j - \theta, 0),
 \qquad \theta \text{ chọn sao cho } \textstyle\sum_j \max(y_j-\theta,0) = n,
\]
thu được từ điều kiện Karush--Kuhn--Tucker của
$\min_{x\in\simplex}\tfrac12\|x-y\|^2$: nó trừ đi một ngưỡng duy nhất rồi
cắt, bảo toàn hiệu giữa các thành phần còn sống. Theo
Mệnh đề~\ref{md:lipschitz} nó còn không giãn. Một phép nhân lại thì bảo toàn
tỉ lệ, và do đó là phép sửa \emph{sai} đối với số lượng.

\subsection{Một điều kiện đủ yếu hơn Giả thiết~\ref{gt:doixung}}

\begin{menhde}[Điều kiện đủ yếu hơn]
\label{md:dieukienyeu}
Giả sử $r_j = 1$ trên các lớp có mặt, có ít nhất một lớp vắng mặt, và khối
lượng lỗi chảy vào thỏa
\begin{equation}
 (Q^Ta)_j = \frac{\sum_k a_k - a_j}{C-1} \qquad \text{với mọi } j .
 \label{eq:luong}
\end{equation}
Khi đó~\eqref{eq:dangdong} đúng chính xác. Ma trận $Q$ đều
kéo theo~\eqref{eq:luong}, nhưng~\eqref{eq:luong} \emph{yếu hơn hẳn}: nó ràng
buộc vector $Q^Ta$ chứ không ràng buộc từng phần tử của $Q$.
\end{menhde}

\begin{proof}
Thế~\eqref{eq:luong} vào $G = -a + Q^Ta$ ta thu lại đúng biểu thức của $G_j$
trong Định lý~\ref{dl:dangdong}, từ đó phần còn lại của chứng minh đó áp dụng
nguyên văn. Để thấy~\eqref{eq:luong} yếu hơn hẳn, lấy $Q$ đều rồi đổi tùy ý
các hàng ứng với lớp vắng mặt: các hàng đó nhân với $a_j = 0$ nên $Q^Ta$ không
đổi, trong khi $Q$ không còn đều.
\end{proof}

Một ví dụ ít tầm thường hơn: khi ba lớp có mặt $0,1,2$ có cùng khối lượng
lỗi, cộng vào $Q$ đều một \emph{vòng luân chuyển} gồm $+\eta$ ở các vị trí
$(0,1),(1,2),(2,0)$ và $-\eta$ ở $(0,2),(1,0),(2,1)$. Các hàng vẫn có tổng
một, và vì ba lớp đó có cùng $a_j$ nên $Q^Ta$ không đổi;
Mục~\ref{sec:kiemchung} kiểm chứng rằng dạng đóng vẫn khôi phục chính xác.

Ngược lại, \eqref{eq:luong} là điều kiện \emph{đủ} chứ không \emph{cần}. Lấy
$C=10$ với các lớp có mặt mang số lượng bằng nhau, và $Q$ gồm các cặp dễ lẫn
$(0,1),(2,3),\ldots$, mỗi lớp gửi một tỉ lệ $s$ khối lượng lỗi sang lớp bạn
của nó. Nếu hai lớp vắng mặt tạo thành một cặp nguyên vẹn thì $G$ chỉ nhận
hai giá trị --- một cho mọi lớp có mặt, một cho mọi lớp vắng mặt --- và dạng
đóng khôi phục chính xác với mọi $s$, \emph{dù}~\eqref{eq:luong} không được
thỏa (phần dư tương đối $0{,}106$ khi $s=0{,}3$ và $0{,}441$ khi $s=0{,}9$).
Ở đây sự chính xác đến từ mục tiêu suy biến (mọi lớp có mặt bằng nhau), chứ
không từ~\eqref{eq:luong}. Cách đọc thực dụng: vùng chính xác của dạng đóng
là một \emph{tập} chứa điểm ``$Q$ đều'', và phần dư của~\eqref{eq:luong} dự
báo tốt mức lệch (Mục~\ref{sec:kiemchung}), nhưng không phải là biên chính
xác của vùng đó.

\subsection{Khi biết hình học: trần lý thuyết}

\begin{giathiet}[Hình học lỗi đã biết và bất khả quy]
\label{gt:hinhhoc}
$Q$ đã biết, row-stochastic, không âm, đường chéo bằng không và bất khả quy;
$r$ đã biết và dương nghiêm ngặt trên các lớp có mặt; có ít nhất một lớp vắng
mặt; $\bar\epsilon>0$.
\end{giathiet}

\begin{dinhly}[Khôi phục khi biết $(r,Q)$]
\label{dl:tongquat}
Dưới Giả thiết~\ref{gt:hinhhoc}, gọi $v>0$ là vector Perron đã chuẩn hóa với
$Q^Tv=v$, $\one^Tv=1$. Mọi nghiệm của~\eqref{eq:hept} đều có dạng
$a = a_0 + tv$ với $a_0 = (I-Q^T)^{+}(-G)$, và các yêu cầu $a \ge 0$ cùng với
sự tồn tại một lớp vắng mặt chọn ra
\begin{equation}
 t = \max_j\Bigl(-\frac{a_{0,j}}{v_j}\Bigr),
 \qquad
 \widehat N_j = n\,\frac{a_j/r_j}{\sum_k a_k/r_k}.
 \label{eq:tongquat}
\end{equation}
Với $G$, $r$, $Q$ chính xác, \eqref{eq:tongquat} bằng đúng số lượng thật.
\end{dinhly}

\begin{proof}
Vì $Q$ row-stochastic và bất khả quy, $\ker(I-Q) = \mathrm{span}\{\one\}$,
nên $\mathrm{range}(I-Q^T) = \one^\perp$; theo~\eqref{eq:tongkhong} ta có
$-G \in \one^\perp$, nên hệ~\eqref{eq:hept} có nghiệm và $a_0$ là một nghiệm.
Theo lý thuyết Perron--Frobenius, không gian hạch của $I-Q^T$ có chiều một,
sinh bởi một $v$ dương nghiêm ngặt, nên mọi nghiệm sai khác nhau một lượng
$tv$. Các ràng buộc $a_{0,j}+tv_j\ge0$
tương đương $t \ge \max_j(-a_{0,j}/v_j)$, và tại $t$ khả thi nhỏ nhất thì ít
nhất một thành phần bằng không. Nghiệm $a$ thật là không âm và có một thành
phần bằng không vì có lớp vắng mặt, nên nó phải là nghiệm biên này. Với các
lớp có mặt, $a_j = N_jr_j\bar\epsilon$, do đó $a_j/r_j = N_j\bar\epsilon$, và
chuẩn hóa các giá trị này về tổng $n$ cho đúng số lượng. Các lớp vắng mặt
đóng góp bằng không, bất kể giá trị dương nào được dùng thế cho $r_j$ của
chúng.
\end{proof}

\begin{hequa}[Biên khi độ khó bị gán sai]
\label{hq:kappa}
Giả sử $Q$ đã biết và nghiệm biên $a$ chính xác, nhưng ước lượng thay bằng
$r_j \equiv 1$. Với $\kappa = r_{\max}/r_{\min}$ trên các lớp có mặt, tỉ lệ
khôi phục theo mẫu $\sum_j \min(\widehat\pi_j, \pi_j)$ ít nhất bằng
$1/\kappa$.
\end{hequa}

\begin{proof}
Ước lượng là $q_j = a_j/A$ với $A=\sum_k a_k$, còn tỉ lệ thật là
$\pi_j = (a_j/r_j)/Z$ với $Z = \sum_k a_k/r_k$. Do đó
$q_j/\pi_j = r_j Z/A$, và $Z/A$ là trung bình có trọng của các $1/r_k$, nên
$Z/A \ge 1/r_{\max}$. Vậy $q_j \ge \pi_j/\kappa$, suy ra
$\min(q_j,\pi_j) \ge \pi_j/\kappa$; cộng lại ta được điều phải chứng minh.
\end{proof}

Định lý~\ref{dl:tongquat} \emph{không} phải một tấn công: $r$ và $Q$ là đầu
vào oracle. Vai trò của nó là cung cấp một trần lý thuyết để đo
Định lý~\ref{dl:dangdong}, với Hệ quả~\ref{hq:kappa} chặn khoảng cách giữa
hai bên. Ba kết quả hợp thành một thang bậc: dạng đóng là ước lượng triển
khai được, dạng tổng quát là trần, và hệ quả xác nhận rằng khoảng cách giữa
chúng suy giảm \emph{có kiểm soát} chứ không sụp đổ.

% ============================================================================
\begin{landscape}
\section{Bảng tóm tắt}
\label{sec:bang}

\footnotesize
\noindent
\begin{tabular}{@{}l l l p{3.8cm} l p{3.8cm} c@{}}
\toprule
\textbf{Phương pháp} & \textbf{Tín hiệu} & \textbf{Vay từ aux} &
\textbf{Cơ chế suy giảm} & \textbf{Căn cứ} &
\textbf{Kết cục khi $\bar\epsilon_{\mathrm{tr}}\to0$} & \textbf{Bất biến} \\
\midrule
iLRG & $\nabla W_j/\nabla b_j$ & không &
Chia cho mẫu số tiêu biến; hệ giả nghịch đảo suy biến &
Bổ đề~\ref{bd:dongnhat} &
Nhiễu số học (phải chặn số mũ thủ công) & --- \\

LLG & $\sum_d \Delta W$ & không &
Không có. Nhưng giả định $r_j\equiv1$, $Q$ đều, và nhân thô $\times m$ &
--- & Vẫn hoạt động & Có \\

LLG$^+$ & $\sum_d \Delta W$ & offset, impact &
Chỉ mẫu dấu $h_1$ còn lại; phần còn lại đóng băng theo offset &
MĐ~\ref{md:llgp} (+ MĐ~\ref{md:khoangcach}) &
Hàm của $h_1$: $\le C$ bit; mô phỏng chỉ $4$ đầu ra, irec dưới dự đoán đều & Một phần \\

ZLG$^+$ & $\sum_d \Delta W$ & $\bar p$, $\bar O$ &
$\widehat N\propto[n\bar p-\tilde G]_+$: khi $\tilde G\to0$ chỉ còn $\bar p$ &
ĐL~\ref{dl:supdo}, HQ~\ref{hq:apdung} &
$n\bar p/\|\bar p\|_1\approx$ phân phối đều. Chồng lấp $0{,}54$ ($\alpha{=}0{,}5$); $0{,}29$ ($\alpha{=}0{,}1$) & Không \\

RLU & $\Delta b$, $\Delta W$ & $\mathrm{shift}$ &
Ba lỗi cộng dồn: sụp đổ phụ trợ; Jensen $e^{\sigma^2/2}$; cắt + chuẩn hóa nhân &
ĐL~\ref{dl:supdo}, HQ~\ref{hq:apdung}, MĐ~\ref{md:jensen},~\ref{md:lipschitz} &
$n\cdot\text{normalize}(\mathrm{shift})$ (lần gọi đầu), độc lập hoàn toàn với nhãn & Không \\
\midrule
VLI (gốc) & $\Delta b$ hoặc $\sum_d\Delta W$ & không &
$f_2,f_3$ \emph{phụ thuộc độ lớn} $G$ & --- &
Xấp xỉ đúng, nhưng 5 siêu tham số phải dò lại & Gần \\

\textbf{VLI}$^*$ \textbf{dạng đóng} & $\Delta b$ & \textbf{không} &
Giả thiết $Q$ đều (MĐ~\ref{md:hoanvi}: lựa chọn \emph{duy nhất}); cần $\ge1$ lớp vắng mặt &
ĐL~\ref{dl:dangdong} &
\textbf{Không đổi.} $\bar\epsilon$ giảm 7 bậc, kết quả y nguyên & \textbf{Có} \\

VLI$^*$ tổng quát & $\Delta b$ & chỉ $(r,Q)$ \emph{đã chuẩn hóa} &
$(r,Q)$ hiện là oracle & ĐL~\ref{dl:tongquat}, HQ~\ref{hq:kappa} &
\textbf{Không đổi} (trần lý thuyết) & Có \\
\bottomrule
\end{tabular}

\vspace{1.4em}
\noindent\normalsize\textbf{Vì sao chỉ VLI$^*$ thoát được}

\vspace{0.5em}
\footnotesize
\noindent
\begin{tabular}{@{}p{4.6cm} p{3.6cm} p{12.4cm}@{}}
\toprule
\textbf{Rào cản} & \textbf{Chặn ai} & \textbf{VLI$^*$ vượt bằng cách nào} \\
\midrule
\textbf{Bổ đề~\ref{bd:chuan}:} $\|G\|_1\le 2nL_{\text{tr}}$ --- toàn bộ tín hiệu co về 0 &
Mọi phương pháp &
Chỉ dùng \emph{tỉ số} $G_j/G_{\max}$; nhân $G$ với hằng số nào cũng không đổi kết quả (MĐ~\ref{md:batbien}) \\

\textbf{ĐL~\ref{dl:supdo}:} $\widehat\pi\to\Phi(c)$, vector độc lập với nhãn &
ZLG$^+$, RLU &
Không tồn tại $c$ ngoài. Thang độ $\bar\epsilon$ được đọc từ $G_{\max}$ --- nhiệt kế nằm ngay trong cửa sổ \\

\textbf{MĐ~\ref{md:khoangcach}:} $\tau_{\text{err}}/\bar\epsilon_{\text{tr}}\to\infty$ (điều kiện đủ) &
LLG$^+$ (ZLG$^+$/RLU không cần) &
Không tham chiếu độ lớn nào từ aux. Riêng $(r,Q)$ tùy chọn đều \emph{vô đơn vị}: $\bar\epsilon$ triệt tiêu theo định nghĩa \\

\textbf{MĐ~\ref{md:llgp}:} đầu ra chỉ còn là hàm của mẫu dấu &
LLG$^+$ &
Không có pha tham lam hay offset; một biểu thức đóng đọc toàn bộ tín hiệu theo tỉ số \\

\textbf{MĐ~\ref{md:jensen}:} Jensen $e^{\sigma^2/2}$, lớn nhất đúng lúc tín hiệu yếu nhất &
RLU & Không lấy softmax của logit trung bình \\

\textbf{MĐ~\ref{md:lipschitz}:} cắt + chia-cho-tổng không Lipschitz (đo được $44\,437\times$) &
RLU, ZLG$^+$ &
Không cần cắt ($\mathrm{raw}\ge0$ tự động) và không cần chuẩn hóa ($\sum\mathrm{raw}=n$ miễn phí, do~\eqref{eq:tongkhong}). Phép chiếu simplex phòng vệ là \emph{cộng} và không giãn \\

Mẫu số $\nabla b\to0$ & iLRG & Gradient bias không bao giờ dùng làm số chia \\

Siêu tham số dò tay & VLI gốc &
$\mathrm{boost}^\star=(C-1)G_{\max}/\sum_j\pos{-G_j}$ được \emph{suy ra}, thay cả 5 hằng số \\
\bottomrule
\end{tabular}

\vspace{1.4em}
\noindent\normalsize\textbf{Sự đánh đổi --- luận điểm trung tâm}

\vspace{0.5em}
\footnotesize
\noindent
\begin{tabular}{@{}l p{6.0cm} p{3.8cm} p{7.8cm}@{}}
\toprule
 & \textbf{Mô hình hình học lỗi} & \textbf{Thang độ} & \textbf{Kết cục khi $\bar\epsilon\to0$} \\
\midrule
RLU, ZLG$^+$, LLG$^+$ &
\emph{Tốt hơn} --- đo $\bar p$, offset, $\mu$ thật từ dữ liệu &
\emph{Phụ thuộc} một độ lớn đo trên phụ trợ &
Sụp về vector độc lập nhãn (ĐL~\ref{dl:supdo}); LLG$^+$ còn $\le C$ bit (MĐ~\ref{md:llgp}) \\
VLI$^*$ &
\emph{Thô hơn} --- giả thiết $Q$ đều, $r_j\equiv1$ &
\emph{Bất biến} &
Sai số \emph{có biên}: $0{,}53$--$1{,}00$ tùy chế độ (HQ~\ref{hq:kappa}, MĐ~\ref{md:dieukienyeu}) \\
\bottomrule
\end{tabular}
\end{landscape}

\noindent
Luận điểm \emph{không} phải ``$Q$ đều là đúng''. Luận điểm là:
\begin{quote}
một mô hình \textbf{sai nhưng bất biến thang độ} thắng một mô hình
\textbf{tốt hơn nhưng phụ thuộc thang độ}, vì sai số hình học là
\emph{có biên} (Hệ quả~\ref{hq:kappa}, Mệnh đề~\ref{md:dieukienyeu}) còn sụp
đổ thang độ là \emph{toàn phần} (Định lý~\ref{dl:supdo}).
\end{quote}
Đây đúng là cấu trúc lập luận giải thích Naive Bayes: giả thiết độc lập có
điều kiện gần như luôn bị vi phạm, nhưng bộ phân lớp vẫn tốt vì quy tắc quyết
định chỉ phụ thuộc một \emph{thứ tự}, không phụ thuộc xác suất được hiệu
chuẩn. Ở đây, theo Mệnh đề~\ref{md:tinhchat}(iii), ước lượng chỉ phụ thuộc
\emph{hiệu} của $\mathrm{raw}$, nên một $Q$ sai là vô hại miễn quan hệ affine
giữa $G$ và $N$ được bảo toàn.

% ============================================================================
\section{Kiểm chứng số}
\label{sec:kiemchung}

\paragraph{Các bảng này là gì.}
Mọi con số dưới đây được tính trên mô hình sinh do Bổ đề~\ref{bd:dongnhat}
hàm ý: cho trước số lượng $N$, vector độ khó $r$, ma trận nhầm lẫn $Q$ và một
thang độ $\bar\epsilon$, ta lập $a_j = N_jr_j\bar\epsilon$, đặt
$G = -(I-Q^T)a$ chính xác, rồi truyền $G$ vào cài đặt. Vì vậy chúng kiểm
chứng phần đại số của Định lý~\ref{dl:dangdong} và~\ref{dl:tongquat}, và đo
độ nhạy khi Giả thiết~\ref{gt:doixung} bị vi phạm một cách có kiểm soát.
\textbf{Chúng không phải thí nghiệm trên một mạng đã huấn luyện} và không
được đọc như vậy; một nghiên cứu thực nghiệm cần các checkpoint ở nhiều mức
độ chính xác và là công việc tiếp theo (Mục~\ref{sec:gioihan}). Với các
phương pháp cơ sở, mô phỏng bám đúng phép chia tỉ lệ của mã nguồn: ZLG$^+$
nhận $v = n\bar p - \tilde G$ (với mọi $w_i = 1$), còn LLG$^+$ nhận tín hiệu
$G/n$, với $\mathrm{offset}_j \approx \bar\epsilon_{\mathrm{aux}}/(C-1)$ và
$\mathrm{impact} = -\bar\epsilon_{\mathrm{aux}}(1+1/C)/B < 0$ như trong mã
nguồn, lấy $\bar\epsilon_{\mathrm{aux}} = 0{,}25$.

\paragraph{Tính chính xác và bất biến thang độ.}
Với $C=4$, $N=(6,4,2,0)$, $Q$ đều, $r\equiv1$, số lượng khôi phục là
$(6,4,2,0)$ chính xác, và vẫn như vậy khi $\bar\epsilon$ --- do đó toàn bộ
tín hiệu --- bị đẩy về không qua bảy bậc độ lớn, xác nhận
Mệnh đề~\ref{md:batbien}.

\begin{center}
\footnotesize
\begin{tabular}{@{}lll@{}}
\toprule
$\bar\epsilon$ & $\|G\|_1$ & Số lượng khôi phục \\
\midrule
$5\times10^{-1}$ & $5{,}33\times10^{0}$  & $(6,4,2,0)$ \\
$1\times10^{-1}$ & $1{,}07\times10^{0}$  & $(6,4,2,0)$ \\
$1\times10^{-3}$ & $1{,}07\times10^{-2}$ & $(6,4,2,0)$ \\
$1\times10^{-8}$ & $1{,}07\times10^{-7}$ & $(6,4,2,0)$ \\
\bottomrule
\end{tabular}
\end{center}

\paragraph{Mệnh đề~\ref{md:lipschitz} đo được.}
Đẩy $\delta \to 0$ dọc theo họ điểm trong chứng minh và đo đạo hàm theo hướng
của $R$, đối chiếu với công thức $4n/(9\delta)$ ($n = 400$):

\begin{center}
\footnotesize
\begin{tabular}{@{}lrr@{}}
\toprule
$\delta$ & tỉ số $L_1$ đo được & $4n/(9\delta)$ \\
\midrule
$10^{-1}$ & $1\,778$ & $1\,778$ \\
$10^{-3}$ & $177\,806$ & $177\,778$ \\
$10^{-5}$ & $17\,783\,464$ & $17\,777\,778$ \\
$10^{-7}$ & $1\,777\,777\,925$ & $1\,777\,777\,778$ \\
\bottomrule
\end{tabular}
\end{center}

Khớp tới ba chữ số có nghĩa trên bốn bậc độ lớn, và phân kỳ như $1/\delta$.
Đối chiếu, trên $3\times10^5$ cặp điểm ngẫu nhiên, hệ số lớn nhất của
$\Pi_{\simplex}$ đo được là $1{,}0000$ trong chuẩn $L_2$ và $1{,}3087$ trong
chuẩn $L_1$ --- cả hai đều dưới biên lý thuyết ($1$ và $\sqrt{C}=3{,}162$).
Đây là lí do Nhận xét~\ref{nx:khongtrichso} từ chối trích một con số đơn lẻ
cho $R$.

\paragraph{Độ nhạy nhãn: cả ZLG$^+$ lẫn LLG$^+$ đều sụp.}
Để tách ``sụp đổ'' khỏi ``suy giảm'', ta cho \emph{hai} vector nhãn khác nhau
tối đa, $N_1=(200,120,41,19,10,5,5,0,0,0)$ và $N_2$ là $N_1$ đảo ngược, với
$\|N_1-N_2\|_1/n = 1{,}900$, rồi đo đầu ra lệch nhau bao nhiêu. Nếu đại lượng
đó $\to 0$ thì ước lượng đã \emph{mù nhãn}: đổi toàn bộ nhãn cũng cho cùng
đáp án. (Ta tránh $N_j = n/C$ đúng, vì khi đó $G_j = 0$ và dấu của nó --- tức
$h_1$ --- do nhiễu dấu phẩy động quyết định; đây chính là mép trong
Mệnh đề~\ref{md:llgp}(i).)

\begin{center}
\footnotesize
\begin{tabular}{@{}lrrrrr@{}}
\toprule
& \multicolumn{3}{c}{LLG$^+$} & \multicolumn{2}{c}{ZLG$^+$} \\
\cmidrule(lr){2-4}\cmidrule(lr){5-6}
$\bar\epsilon_{\mathrm{tr}}$ & $\|o_1-o_2\|_1/n$ & cAcc & irec
 & $\|o_1-o_2\|_1/n$ & irec \\
\midrule
$10^{-1}$ & $0{,}700$ & $0{,}50$ & $0{,}388$ & $0{,}040$ & $0{,}405$ \\
$10^{-2}$ & $0{,}050$ & $0{,}60$ & $0{,}200$ & $0{,}000$ & $0{,}398$ \\
$10^{-3}$ & $0{,}050$ & $0{,}60$ & $0{,}200$ & $0{,}000$ & $0{,}398$ \\
$10^{-4}$ & $0{,}050$ & $0{,}60$ & $0{,}200$ & $0{,}000$ & $0{,}398$ \\
$10^{-6}$ & $0{,}050$ & $0{,}60$ & $0{,}200$ & $0{,}000$ & $0{,}398$ \\
$10^{-8}$ & $\mathbf{0{,}050}$ & $0{,}60$ & $\mathbf{0{,}200}$ & $\mathbf{0{,}000}$ & $\mathbf{0{,}398}$ \\
\bottomrule
\end{tabular}
\end{center}

Hai kết luận. (a) ZLG$^+$ mất hoàn toàn độ nhạy nhãn, và irec của nó dừng ở
$0{,}398$ --- \emph{đúng bằng} $\sum_j\min(n/C, N_{1,j})/n = 159/400$, tức
giá trị của việc dự đoán phân phối đều, đúng như Hệ quả~\ref{hq:apdung}.
(b) LLG$^+$ \emph{cũng} gần như mù nhãn: hai bộ nhãn đối nghịch chỉ khác nhau
ở $h_1$ --- $\{0,1,2\}$ so với $\{7,8,9\}$ --- và điều đó chỉ đổi $2$ trong
$40$ suất, tức $\|o_1-o_2\|_1/n = 0{,}05$. Đầu ra khớp $F(h_1)$ của
Mệnh đề~\ref{md:llgp} trong $500/500$ trường hợp ngẫu nhiên, và trên $3000$
vector nhãn ngẫu nhiên chỉ có $4$ đầu ra khác nhau. Vì phần lớn suất được
phân theo offset chứ không theo nhãn, irec của LLG$^+$ ($0{,}200$) còn
\emph{thấp hơn} dự đoán đều ($0{,}398$).

\paragraph{Biên lượng tử hóa của Hệ quả~\ref{hq:apdung}(ii).}
Cùng cấu hình ZLG$^+$, nhưng đo trực tiếp
$\|\widehat N - R(c)\|_1$ trên $4000$ vector nhãn ngẫu nhiên
mỗi mức $\bar\epsilon$ (Dirichlet với $\alpha \in \{0{,}1;\,0{,}4;\,1;\,5\}$).
Ta lấy $\bar p_j = 0{,}7512$ để $c_j = 300{,}48$ nằm xa mép của \texttt{int};
cột ``trên mép'' dùng $\bar p$ được dựng cố ý sao cho $R(c)_0 = 40{,}5$ nằm
đúng trên mép của \texttt{round}.

\begin{center}
\footnotesize
\begin{tabular}{@{}lrrrr@{}}
\toprule
$\bar\epsilon_{\mathrm{tr}}$ & $L_\Phi\beta\|\Gamma\|_1$ & biên~\eqref{eq:bienluongtu}
 & $\|\widehat N - R(c)\|_1$ ngoài mép & trên mép \\
\midrule
$10^{-2}$ & $2{,}130$ & $9{,}792$ & $\mathbf{0}$ & $\mathbf{1}$ \\
$10^{-3}$ & $0{,}213$ & $7{,}875$ & $\mathbf{0}$ & $\mathbf{1}$ \\
$10^{-4}$ & $0{,}021$ & $7{,}684$ & $\mathbf{0}$ & $\mathbf{1}$ \\
$10^{-6}$ & $0{,}000$ & $7{,}663$ & $\mathbf{0}$ & $\mathbf{1}$ \\
$10^{-8}$ & $0{,}000$ & $7{,}662$ & $\mathbf{0}$ & $\mathbf{1}$ \\
\bottomrule
\end{tabular}
\end{center}

Ở đây $\delta = n\sum_j \bar p_j \approx 3005$, nên $L_\Phi = 2n/\delta
\approx 0{,}266$ và $\gamma_{\mathcal{Q}}C = (L_\Phi+\tfrac12)C \approx
7{,}66$. Biên không bị vượt, và thực tế còn lỏng hơn nhiều: $0$ khi $R(c)$
nằm ngoài mép (đóng băng chính xác như Nhận xét~\ref{nx:dongbang} tiên đoán),
và đúng \emph{một} count khi nằm trên mép --- tức $1/400 = 0{,}25\%$. Ngưỡng
đủ của Nhận xét~\ref{nx:dongbang} là
$\bar\epsilon_0 = \rho/(2\beta nL_\Phi) \approx 2{,}3\times10^{-3}$, trong khi
mô phỏng đã đóng băng từ $10^{-2}$: ngưỡng đó bảo thủ. Trên $\mathrm{irec}$,
\eqref{eq:bienirec} cho biên $1{,}92\%$ còn độ lệch quan sát được là $0$ hoặc
$0{,}25\%$.

\paragraph{Đối chứng cho Mệnh đề~\ref{md:conthongtin}: thông tin vẫn còn.}
Cùng cấu hình ZLG$^+$ và cùng nhãn thật $N_1$ như trên.
Cột cuối lấy $e = R(v)-R(c)$ --- một đại lượng kẻ tấn công \emph{tính được},
vì $c$ do chính họ đo --- rồi đọc $-e$ bằng dạng đóng không thang độ
\eqref{eq:dangdong} thay vì bằng quy tắc của ZLG$^+$.

\begin{center}
\footnotesize
\begin{tabular}{@{}lrrr@{}}
\toprule
$\bar\epsilon_{\mathrm{tr}}$ & $\|e\|_\infty$ & irec của ZLG$^+$
 & irec khi đọc $-e$ theo tỉ số \\
\midrule
$10^{-2}$ & $2{,}37\times10^{-1}$ & $0{,}398$ & $\mathbf{1{,}000}$ \\
$10^{-4}$ & $2{,}37\times10^{-3}$ & $0{,}398$ & $\mathbf{1{,}000}$ \\
$10^{-6}$ & $2{,}37\times10^{-5}$ & $0{,}398$ & $\mathbf{1{,}000}$ \\
$10^{-8}$ & $2{,}37\times10^{-7}$ & $0{,}398$ & $\mathbf{1{,}000}$ \\
\bottomrule
\end{tabular}
\end{center}

Hai cách đọc \emph{cùng một} $R(v)$: quy tắc của ZLG$^+$ cho $0{,}398$ ---
đúng giá trị của dự đoán đều --- còn phép đọc theo tỉ số khôi phục
\emph{chính xác tuyệt đối}, đúng như Mệnh đề~\ref{md:conthongtin} dự đoán
($-e$ tỉ lệ với $\Gamma$). Vậy việc tín hiệu co nhỏ \emph{không} xoá thông
tin; nó chỉ làm thông tin trở nên nhỏ, và chính thiết kế phụ thuộc thang độ
mới đánh mất nó (Mục~\ref{sec:tuongtac}).

\paragraph{Hình dạng so với độ lớn: $Q$ chuyển giao, nhưng không cố định.}
Dùng mô hình logit $z = s\cdot\mathrm{diff}[y]\cdot M[y] + \sigma\varepsilon$,
với train $\sigma=1{,}0$ và phụ trợ $\sigma=2{,}5$. Tăng $s$ mô phỏng mô hình
tốt dần. Cột ``$Q$ trôi'' so $Q(s)$ với $Q(s{=}1)$ (\emph{qua thời gian});
hai cột cuối so phụ trợ với train \emph{tại cùng một $s$}.

\begin{center}
\footnotesize
\begin{tabular}{@{}lrrrrr@{}}
\toprule
$s$ & $\bar\epsilon_{\mathrm{tr}}$ & $Q$ trôi theo $s$ & độ tập trung $Q$
 & $Q$ lệch aux/train & offset lệch aux/train \\
\midrule
$1{,}0$ & $7{,}9\times10^{-1}$ & $0{,}000$ & $0{,}142$ & $0{,}094$ & $0{,}1$ \\
$2{,}0$ & $6{,}3\times10^{-1}$ & $0{,}165$ & $0{,}178$ & $0{,}149$ & $0{,}2$ \\
$3{,}0$ & $4{,}4\times10^{-1}$ & $0{,}332$ & $0{,}219$ & $0{,}215$ & $0{,}5$ \\
$4{,}0$ & $2{,}8\times10^{-1}$ & $0{,}608$ & $0{,}285$ & $0{,}292$ & $0{,}9$ \\
$6{,}0$ & $9{,}7\times10^{-2}$ & $1{,}130$ & $0{,}411$ & $0{,}350$ & $2{,}3$ \\
$8{,}0$ & $3{,}3\times10^{-2}$ & $\mathbf{1{,}515}$ & $\mathbf{0{,}510}$ & $\mathbf{0{,}302}$ & $\mathbf{3{,}5}$ \\
\bottomrule
\end{tabular}
\end{center}

Ba điều cần đọc đúng. (a) $Q$ \emph{không} cố định qua huấn luyện: nó trôi
$151\%$. Tính bất biến ở Mục~\ref{sec:kihieu} là bất biến \emph{đối với một
hệ số chung}, không phải tính ổn định theo thời gian. (b) Nhưng cái mà tấn
công cần là cột ``$Q$ lệch aux/train'', và nó \emph{bão hòa} quanh $0{,}30$,
trong khi ``offset lệch'' \emph{phân kỳ} tới $3{,}5$ ($350\%$) và chưa dừng.
Đây là khác biệt về bản chất, và là bằng chứng cho giả thuyết ở
Mục~\ref{sec:gioihan}. (c) Cột ``độ tập trung'' (trung bình của
$\max_k Q_{jk}$; giá trị đều là $1/(C-1)\approx0{,}111$) tăng từ $0{,}142$
lên $0{,}510$: \emph{càng huấn luyện, $Q$ càng rời xa ma trận đều}. Đối chiếu
bảng nhầm lẫn có cấu trúc dưới đây, $0{,}510$ ứng với mức irec khoảng
$0{,}88$. Đây là một giới hạn, nêu ở Mục~\ref{sec:gioihan}.

\paragraph{Nhầm lẫn có cấu trúc.}
$C=10$, $N=(120,90,70,50,30,20,10,5,0,0)$, $r\equiv1$, và $Q$ dựng từ các cặp
dễ lẫn với cường độ $s$ (giá trị $s=1/(C-1)\approx0{,}111$ là trường hợp
đều). Đây là vi phạm đáng lo nhất, vì mô hình thực luôn lẫn các cặp cụ thể.

\begin{center}
\footnotesize
\begin{tabular}{@{}lr@{}}
\toprule
Cường độ cặp $s$ & Khôi phục theo mẫu \\
\midrule
$0{,}111$ (đều) & $1{,}000$ \\
$0{,}200$ & $0{,}982$ \\
$0{,}300$ & $0{,}962$ \\
$0{,}500$ & $0{,}884$ \\
$0{,}700$ & $0{,}734$ \\
$0{,}900$ & $0{,}638$ \\
\bottomrule
\end{tabular}
\end{center}

\paragraph{Chính xác bên trong và bên ngoài điều kiện~\eqref{eq:luong}.}
Trước hết, ví dụ vòng luân chuyển của Mệnh đề~\ref{md:dieukienyeu} (tám lớp
có mặt, mỗi lớp $70$ mẫu, $\eta = 0{,}05$) cho một $Q$ không đều nhưng thỏa
\eqref{eq:luong}, và dạng đóng khôi phục đúng $(70,\ldots,70,0,0)$. Thứ hai,
với $Q$ cấu trúc cặp, các lớp có mặt mang số lượng bằng nhau và hai lớp vắng
mặt tạo thành một cặp nguyên vẹn, dạng đóng cũng khôi phục chính xác với mọi
cường độ đã thử --- dù ở đây \eqref{eq:luong} \emph{không} được thỏa. Dòng đầu
bảng dưới đây là trường hợp đó; nó cho thấy \eqref{eq:luong} là đủ chứ không
cần. Ba dòng sau phá sự đối xứng của mục tiêu và cho thấy mức lệch quay lại.

\begin{center}
\footnotesize
\begin{tabular}{@{}lrr@{}}
\toprule
Cấu hình & $s=0{,}3$ & $s=0{,}9$ \\
\midrule
Có mặt cân bằng, cặp vắng mặt nguyên vẹn & $1{,}000$ & $1{,}000$ \\
Có mặt cân bằng, cặp bị chia giữa có mặt/vắng mặt & $0{,}964$ & $0{,}864$ \\
Mất cân bằng nhẹ, cặp nguyên vẹn & $0{,}985$ & $0{,}667$ \\
Mất cân bằng mạnh, cặp nguyên vẹn & $0{,}940$ & $0{,}530$ \\
\bottomrule
\end{tabular}
\end{center}

\paragraph{Phần dư của~\eqref{eq:luong} tiên đoán mức mất mát.}
Trên $400$ cấu hình ngẫu nhiên (thay đổi $N$, số lớp có mặt, và
$Q \in \{\text{đều}, \text{cặp}, \text{vòng}\}$), đặt
$\mathrm{defect} = \|Q^Ta - (\one^Ta\,\one - a)/(C-1)\|_2/\|a\|_2$. Tương
quan giữa defect và khôi phục là $-0{,}873$, và defect bằng không cho kết quả
chính xác trong \emph{toàn bộ} $138$ trường hợp xuất hiện, phù hợp với tính
\emph{đủ} của Mệnh đề~\ref{md:dieukienyeu}.

\begin{center}
\footnotesize
\begin{tabular}{@{}lrr@{}}
\toprule
Khoảng defect & Số cấu hình & Khôi phục trung bình \\
\midrule
$0$ (chính xác) & $138$ & $1{,}000$ \\
$(0;\,0{,}05)$  & $17$  & $0{,}958$ \\
$[0{,}05;\,0{,}1)$ & $39$ & $0{,}917$ \\
$[0{,}1;\,0{,}3)$ & $65$  & $0{,}832$ \\
$[0{,}3;\,\infty)$ & $141$ & $0{,}716$ \\
\bottomrule
\end{tabular}
\end{center}

\paragraph{Độ phân tán độ khó, đối chiếu Hệ quả~\ref{hq:kappa}.}
$Q$ đều, $r$ phân tán tuyến tính, ước lượng dùng $r\equiv1$. Biên của hệ quả
là đúng nhưng rất lỏng; suy giảm thực tế mượt hơn nhiều bậc.

\begin{center}
\footnotesize
\begin{tabular}{@{}lrr@{}}
\toprule
$\kappa$ & Khôi phục theo mẫu & Biên $1/\kappa$ \\
\midrule
$1{,}00$  & $1{,}000$ & $1{,}000$ \\
$1{,}39$  & $0{,}965$ & $0{,}720$ \\
$1{,}78$  & $0{,}934$ & $0{,}562$ \\
$3{,}33$  & $0{,}853$ & $0{,}300$ \\
$8{,}00$  & $0{,}744$ & $0{,}125$ \\
$39{,}11$ & $0{,}635$ & $0{,}026$ \\
\bottomrule
\end{tabular}
\end{center}

\paragraph{Nhiễu động trên cái mốc.}
Nhiễu Gauss với độ lệch chuẩn $\sigma\cdot\|G\|_\infty$ cộng vào $G$, lặp
$2000$ lần. Cột cuối cho biết nhiễu đẩy $\argmaxop_j G_j$ sang một lớp
\emph{có mặt} bao nhiêu phần trăm thời gian --- điều phá hủy cách hiểu
$G_{\max}$ như một nhiệt kế.

\begin{center}
\footnotesize
\begin{tabular}{@{}lrr@{}}
\toprule
$\sigma$ & Khôi phục theo mẫu & Mốc rơi vào lớp có mặt \\
\midrule
$0\%$  & $1{,}000$ & $0{,}0\%$ \\
$1\%$  & $0{,}991$ & $0{,}0\%$ \\
$5\%$  & $0{,}954$ & $8{,}6\%$ \\
$10\%$ & $0{,}914$ & $22{,}6\%$ \\
$25\%$ & $0{,}820$ & $44{,}0\%$ \\
$50\%$ & $0{,}722$ & $53{,}4\%$ \\
\bottomrule
\end{tabular}
\end{center}

\paragraph{Yêu cầu về lớp vắng mặt.}
Khi không có lớp nào vắng mặt, công thức bị buộc phải gán không cho thành
phần $\argmaxop$ của nó, nên lớp có mặt nhỏ nhất bị hy sinh. Thiệt hại do đó
tỉ lệ với \emph{tỉ trọng} của lớp đó trong $n$, chứ không tỉ lệ với bản thân
việc vắng mặt: với cấu hình $C=10$ đuôi dài ở trên, loại bỏ hoàn toàn các lớp
vắng mặt vẫn cho $0{,}975$, trong khi một cửa sổ nhỏ và gần cân bằng thì
thiệt hại nhiều hơn đáng kể. Một cửa sổ cân bằng hoàn hảo dưới
Giả thiết~\ref{gt:doixung} cho $G \equiv 0$, và cài đặt từ chối trả lời thay
vì bịa ra một đáp án.

% ============================================================================
\section{Giới hạn và các nguy cơ đối với tính hợp lệ}
\label{sec:gioihan}

\begin{enumerate}
\item \textbf{$(r,Q)$ là đầu vào oracle.} Định lý~\ref{dl:tongquat} giả định
  đã biết chúng. Việc $\bar\epsilon$ triệt tiêu khỏi cả hai theo cách xây
  dựng (Mục~\ref{sec:kihieu}) là một lí do để kì vọng chúng chuyển giao từ
  dữ liệu phụ trợ sang dữ liệu huấn luyện tốt hơn các độ lớn tuyệt đối mà
  các phương pháp cũ dùng, nhưng đây là một \emph{giả thuyết có cơ sở, không
  phải một định lý}, và phải được đo. Khoảng cách từ $0{,}638$ tới $1{,}000$
  trong bảng nhầm lẫn có cấu trúc chính là giá trị định lượng của việc giải
  được vấn đề này.

\item \textbf{$\kappa$ nở ra theo độ huấn luyện.} Vì
  $r_j = \bar\epsilon_j/\bar\epsilon$ là một tỉ số, $\bar\epsilon$ triệt
  tiêu, nhưng \emph{độ phân tán} của $r$ thì không nhất thiết ổn định: các
  lớp dễ có thể đạt $\bar\epsilon_j\approx10^{-6}$ trong khi các lớp khó vẫn
  ở $10^{-2}$, cho $\kappa \approx 10^4$, ở nơi một mô hình chưa huấn luyện
  đủ chỉ có $\kappa \approx 2$. Vậy chế độ mà VLI$^*$ thắng cũng chính là chế
  độ mà Giả thiết~\ref{gt:doixung} kém chính xác nhất. Đây là \emph{tension}
  trung tâm của hướng tiếp cận và phải được báo cáo, không được che; nó đo
  được trực tiếp từ các checkpoint có độ chính xác tăng dần.

\item \textbf{$Q$ tập trung dần, cũng theo độ huấn luyện.} Đây là một kênh
  \emph{thứ hai, độc lập} với $\kappa$, ăn vào cùng một chỗ. Số liệu ở
  Mục~\ref{sec:kiemchung} cho độ tập trung của $Q$ (trung bình
  $\max_k Q_{jk}$) tăng từ $0{,}142$ lên $0{,}510$ khi mô hình tốt dần, trong
  khi giá trị của ma trận đều là $1/(C-1)\approx0{,}111$. Nói cách khác, mô
  hình well-trained chỉ còn lẫn các cặp \emph{thật sự} khó, nên một nửa khối
  lượng lỗi dồn vào một lớp duy nhất --- và
  Giả thiết~\ref{gt:doixung} trở nên sai dần đúng vào chế độ mà VLI$^*$ được
  thiết kế cho. Đối chiếu bảng nhầm lẫn có cấu trúc, mức $0{,}510$ ứng với
  irec khoảng $0{,}88$. Cùng với $\kappa$, đây là hai giới hạn quan trọng
  nhất của dạng đóng, và cả hai đều đo được từ cột $Q$ của script chẩn đoán.

\item \textbf{Tiên đoán khả phủ chứng về cơ chế sụp đổ.} Lý thuyết cho hai
  tiên đoán khác nhau theo phương pháp. (a) Với ZLG$^+$ và RLU, sự sụp đổ chỉ
  cần $\bar\epsilon_{\mathrm{tr}}\to0$ (Hệ quả~\ref{hq:apdung}), nên mức suy
  giảm phải theo \emph{độ chính xác huấn luyện}, không theo khoảng cách tổng
  quát hóa; và irec cuối cùng phải xấp xỉ độ chồng lấp giữa phân phối nhãn của
  client với $\Phi(c)$ (gần đều) --- tức phụ thuộc vào $\alpha$ của phân hoạch
  Dirichlet, không phụ thuộc vào tập dữ liệu. (b) Với LLG$^+$, sự đóng băng đòi
  $\bar\epsilon_{\mathrm{tr}} \ll \bar\epsilon_{\mathrm{aux}}$
  (Mệnh đề~\ref{md:llgp}), nên trên một mô hình có khoảng cách tổng quát hóa
  nhỏ (MNIST + LeNet5) LLG$^+$ phải hoạt động lâu hơn so với CIFAR100 +
  ResNet18. Nếu irec cuối của ZLG$^+$ khác xa độ chồng lấp tiên đoán, hoặc
  LLG$^+$ trên MNIST sụp sớm như trên CIFAR100, thì cơ chế được đề xuất ở đây
  \emph{sai} ở phần tương ứng.

\item \textbf{Cái mốc chỉ là một toạ độ.} $G_{\max}$ mang toàn bộ việc hiệu
  chuẩn thang độ, nên nhiễu cộng vừa phải --- loại mà một cập nhật có bảo mật
  vi phân hoặc bị nén sẽ tạo ra --- làm VLI$^*$ suy giảm nhiều hơn so với một
  ước lượng lấy trung bình trên nhiều toạ độ. Điều này gợi ra một hướng phòng
  thủ tự nhiên: nhiễu động chỉ cần đủ để che lớp nào vắng mặt, chứ không cần
  nhiễu bảo mật mạnh.

\item \textbf{Đọc gradient là một tiện nghi hộp trắng.} Các đồng nhất thức
  đòi gradient cross-entropy. Một cài đặt đọc bộ đệm gradient của tham số
  \emph{trước} bước optimizer thu được đúng đại lượng đó, nhưng một server
  thật quan sát một cập nhật trọng số tích lũy, mà dưới momentum, weight
  decay, hoặc nhiều bước cục bộ thì không bằng $\sum_i(p_i-e_{y_i})$.
  Momentum bằng không là \emph{một yêu cầu của phép suy luận}, không phải một
  tùy chọn cho tiện.

\item \textbf{Các số hạng chỉnh proximal và động làm nhiễm tín hiệu.} Số hạng
  FedProx đóng góp $\mu(b-b_g)$ vào gradient bias, vốn không có tổng bằng
  không và cũng không có dạng~\eqref{eq:hept}; khi đó
  cả~\eqref{eq:tongkhong} lẫn cái mốc $G_{\max}$ đều sai. Dạng đóng nên được
  giới hạn cho FedAvg thuần, và việc bản heuristic \emph{có vẻ} vẫn chạy dưới
  FedProx không có cơ sở lý thuyết nào.

\item \textbf{Phép gộp.} Nếu server chỉ quan sát trung bình của nhiều client,
  mọi phương pháp trong báo cáo này khôi phục \emph{hợp} các nhãn của chúng,
  không phải một phân rã theo từng client.

\item \textbf{Phần số là dựa trên mô hình.} Như đã nêu ở
  Mục~\ref{sec:kiemchung}, các bảng đánh giá mô hình sinh, không phải một
  mạng đã huấn luyện. Chương trình thực nghiệm sẽ hoàn thiện báo cáo này là:
  tạo các checkpoint trải một khoảng độ chính xác huấn luyện; đo
  $\bar\epsilon_{\mathrm{tr}}$, $\bar\epsilon_{\mathrm{aux}}$, $\kappa$ và
  $Q$ trên từng checkpoint; kiểm chứng trực tiếp Bổ đề~\ref{bd:chuan} và
  Mệnh đề~\ref{md:khoangcach}; và vẽ tỉ lệ khôi phục theo
  $\bar\epsilon_{\mathrm{tr}}$ cho mọi phương pháp, báo cáo riêng tỉ lệ các
  cửa sổ có chứa lớp vắng mặt.
\end{enumerate}

% ============================================================================
\begin{thebibliography}{9}
\bibitem{ilrg} J.~Ma và cộng sự.
  \emph{Instance-wise Batch Label Restoration via Gradients in Federated
  Learning}. ICLR, 2023.
\bibitem{llg} A.~Wainakh và cộng sự.
  \emph{User-Level Label Leakage from Gradients in Federated Learning}.
  arXiv:2105.09369, 2021.
\bibitem{zlg} J.~Geng và cộng sự.
  \emph{Towards General Deep Leakage in Federated Learning}.
  arXiv:2110.09074, 2021.
\bibitem{rlu} H.~Chen và H.~Vikalo.
  \emph{Recovering Labels from Local Updates in Federated Learning}.
  ICML, 2024.
\bibitem{duchi} J.~Duchi, S.~Shalev-Shwartz, Y.~Singer và T.~Chandra.
  \emph{Efficient Projections onto the $\ell_1$-Ball for Learning in High
  Dimensions}. ICML, 2008.
\bibitem{nb} P.~Domingos và M.~Pazzani.
  \emph{On the Optimality of the Simple Bayesian Classifier under Zero-One
  Loss}. Machine Learning, 29:103--130, 1997.
\end{thebibliography}

\end{document}
